\chapter{不确定性规划：分支场景、鲁棒 Tube MPC 与 CBF 安全滤波}
\label{ch:plan-robust}

% ════════════════════════════════════════════════════════════════
%  不确定性规划：全吸收章。
%  吸收 分支/场景规划（Hardy-Campbell 2013 / MPDM Cunningham 2015 /
%  EUDM 2020 / EPSILON 2021 / MARC 2023 / multi-stage NMPC Lucia 2013）
%  + 鲁棒 Tube MPC（Mayne 2005 / Homothetic 2012 / Elastic 2016 / RPI Rakovic 2005）
%  + 可达安全包络（FaSTrack 2017 / RTD 2020 / Funnel 2017）
%  + GP-MPC（Hewing 2020 / Kabzan 2019 / Torrente 2021）
%  + CBF-QP 安全滤波（Ames 2017/2019 / HOCBF Xiao-Belta 2022 / Grandia 2021）。
%  跨部去重：CBF/CLF、HJ 可达、鲁棒/随机 MPC 的最深推导留待控制部，本章侧重规划/决策用途。
%  教学层遵循 v5.0 理论教学规范（动机→反面→理论→陷阱→练习 + 符号表/速查/故障排查）。
% ════════════════════════════════════════════════════════════════

\begin{practice}[前置自测]
开始之前，先试答下列问题。若已能流畅作答，可只速读本章表格与速查卡；若卡住，正是本章要为你解决的。
\begin{enumerate}
  \item 写出一个标准 MPC（模型预测控制）的有限时域最优控制问题：决策变量、代价、动力学约束、状态/控制约束各是什么？它对“未来”做了什么隐含假设？
  \item 一辆自动驾驶车在路口看到一辆横向来车，它“可能让行也可能抢行”。如果你的规划器只赌“最可能”的那个未来，会有什么风险？
  \item Minkowski 和 $\mathcal{A}\oplus\mathcal{B}$ 是什么？$[-1,1]\oplus[-2,2]$ 等于什么？Pontryagin 差 $\mathcal{A}\ominus\mathcal{B}$ 又是什么？
  \item 离散线性系统 $\mathbf{e}_{k+1}=\mathbf{M}\mathbf{e}_k$ 在什么条件下渐近稳定？谱半径 $\rho(\mathbf{M})$ 在其中起什么作用？
  \item 李雅普诺夫函数 $V(\mathbf{x})$ 用来证明什么？把它的“$\dot V$ 沿轨迹递减”换一个不等号方向，能不能用来论证“状态永远不离开某个安全集”？
\end{enumerate}
\noindent\emph{答案线索}：1、2 是全章的动机地基（\cref{sec:pr-branch-form}）；3 是 Tube MPC 约束收紧的工具（\cref{sec:pr-tube}）；4 是偏差动力学稳定性（\cref{sec:pr-tube}）；5 正是 CBF 安全条件（\cref{sec:pr-cbf}）。MPC、LQR、QP、李雅普诺夫的系统讲解见 \cref{ch:control}（\cref{sec:ctrl-mpc,sec:ctrl-lqr,sec:ctrl-lyapunov}）。
\end{practice}

\section*{本章目标}
学完本章，你应当能够：
\begin{enumerate}
  \item \textbf{辨析}规控中的四类不确定性（认知 / 偶然 / 模型 / 环境），并说清“连续小扰动”与“离散多模态意图”为何需要本质不同的处理；
  \item \textbf{推导}分支 / contingency 规划的基本形式化——场景加权代价 + 共享干（非预期约束）+ 各分支独立优化，并解释三个关键设计旋钮（分支点、分支数、场景概率）；
  \item \textbf{说清} MPDM 的策略枚举 + 闭环前向仿真、EUDM 的 DCP-Tree + CFB 如何把组合爆炸压回实时、EPSILON 的分层架构，以及 MARC 的动态分支点 + CVaR 风险感知；
  \item \textbf{建立}分支规划与多阶段（multi-stage）NMPC 场景树的数学同构，并对比分支规划与 MCTS 的结构差异；
  \item \textbf{推导} Tube MPC 的核心分解（名义系统 + 反馈修正 + RPI 集），独立完成 Rigid Tube 的约束收紧，并在 Rigid / Homothetic / Elastic 三代间选型；
  \item \textbf{解释} FaSTrack / RTD / Funnel 三种可达性方法如何给出“与规划器解耦”的安全包络，并把它们与 RPI 集统一在“鲁棒不变 / 可达集”框架下；
  \item \textbf{搭建} GP-MPC 的数据驱动架构（名义模型 + GP 残差 → 协方差前向传播 → 机会约束收紧），理解“均值修正动力学、方差驱动保守度”；
  \item \textbf{实现} CBF-QP 安全滤波器，处理高阶相对度（HOCBF）与 CLF-CBF 联立，理解它如何以“最小修正”给任意名义控制器兜安全底线、与 MPC 形成层级互补。
\end{enumerate}

\subsection*{本章知识导航}
本章是一条“\emph{当未来不确定时，如何在规划与控制中显式、可证地携带这份不确定}”的主线。不确定性分两副面孔——\textbf{离散多模态}（他人意图：让 / 不让）与\textbf{连续有界 / 概率}（动力学扰动、模型误差），对应两大工具族：\emph{分支 / 场景规划}处理前者，\emph{鲁棒规划与安全滤波}处理后者。结构如下：

\begin{center}
\begin{forest} navtree
[{不确定性分类\\认知/偶然/模型/环境}
  [{离散多模态\\分支/场景规划}
    [{MPDM/EUDM\\EPSILON/MARC}
      [{multi-stage NMPC\\场景树（同构）}]]]
  [{连续有界/概率\\鲁棒规划+安全滤波}
    [{Tube MPC\\Rigid/Homo/Elastic}
      [{可达包络\\FaSTrack/RTD/Funnel}
        [{GP-MPC + CBF-QP\\数据残差 + 安全滤波}]]]]
]
\end{forest}
\end{center}

\noindent\textbf{推荐路径}：\cref{sec:pr-uncertainty}（不确定性分类）是全章地基，给出“该用哪类工具”的判据；\cref{sec:pr-branch-form}（分支规划形式化）与 \cref{sec:pr-tube}（Tube MPC 数学本质）是两条主线各自的核心，务必读透。想动手做交互式驾驶决策 → \cref{sec:pr-mpdm,sec:pr-eudm,sec:pr-epsilon}（MPDM/EUDM/EPSILON）；想在真机上快速获得安全保障 → \cref{sec:pr-cbf}（CBF-QP，工程上最快见效）。带“前沿”标记的 MARC（\cref{sec:pr-marc}）、可达包络（\cref{sec:pr-reach}）、GP-MPC（\cref{sec:pr-gpmpc}）为进阶。

\subsection*{前置知识桥接}
本章建立在 \cref{ch:planning}（运动规划导论）与 \cref{ch:control}（控制导论）之上，这里用几行重新激活，便于不翻回去也能跟上。
\begin{itemize}
  \item \textbf{MPC 的优化结构}（\cref{sec:ctrl-mpc}）：MPC 在每个控制周期解一个有限时域最优控制问题——决策一条控制序列，最小化预测时域上的代价，满足动力学、状态、控制约束，施加首控制后\emph{滚动重解}（receding horizon）。本章把“一条名义轨迹”升级成“一条轨迹 + 一条管道”，或升级成“一棵轨迹树”。
  \item \textbf{LQR 与 Riccati}（\cref{sec:ctrl-lqr}）：Tube MPC 的镇定反馈增益 $\mathbf{K}$ 常由对名义系统做 LQR 得到（离散代数 Riccati 方程的解）；CBF-QP 的“最小修正”本质是一个凸 QP，与 MPC 凝聚出的 QP 同属一类（\cref{sec:ctrl-mpc-qp}）。
  \item \textbf{李雅普诺夫稳定性}（\cref{sec:ctrl-lyapunov,sec:ctrl-lyap-direct}）：偏差动力学的镇定靠 $A+BK$ 稳定（谱半径 $<1$）；CBF 安全条件 $\dot h\ge-\alpha(h)$ 是把李雅普诺夫的 $\dot V\le 0$ “换一个不等号方向”用于前向不变性论证。
  \item \textbf{规划三范式}（\cref{ch:planning}）：搜索式（A$^*$、\cref{sec:plan-astar}）、采样式（RRT、\cref{sec:plan-sampling}）、优化式（\cref{sec:plan-traj}）。本章 §\ref{sec:pr-mcts} 把分支规划与蒙特卡洛树搜索（MCTS）放进同一坐标系对比。
\end{itemize}

\subsection*{如果跳过本章会怎样}
\begin{itemize}
  \item \textbf{僵住的机器人（frozen robot）}：你的机器人在人群 / 车流中导航，预测每个动态障碍的未来后发现“几乎所有方向都可能撞”，单轨迹规划器找不到可行解、停在原地不动。根因是它假设“未来唯一确定”并对最坏预测买单——\cref{sec:pr-branch-form} 的分支规划用“对多个离散未来各留一条出路、当前只走公共安全动作”化解。
  \item \textbf{仿真好、真机摔}：名义 MPC 控四足在仿真里走得稳，真机一上斜坡就摔——摩擦、电机延迟、负载偏差偏离名义模型，而名义 MPC 没有余量吸收偏差。\cref{sec:pr-tube} 的 Tube MPC 给名义轨迹外围预留一条“安全管道”框住偏差。
  \item \textbf{RL 策略不敢上车}：你训了个表现好但偶发危险动作的 RL 策略，没有任何安全保证。\cref{sec:pr-cbf} 的 CBF 安全滤波套在\emph{任意}名义控制器（含黑箱 RL）外面，以“最小修正”实时拦截危险动作，给“不可证明安全的策略”一道形式化底线。
\end{itemize}

\begin{table}[H]\centering\small
\caption{本章阅读方式建议}
\label{tab:pr-reading}
\begin{tabular}{lll}
\toprule
阅读方式 & 时间 & 适合谁 \\
\midrule
精读（含全部推导与练习） & 8--12 小时 & 首次系统学习不确定性规控、要做安全关键系统 \\
速读（抓两条主线本质 + 选型表） & 3--4 小时 & 有 MPC 背景、想对齐本书记号与方法谱系 \\
速查（只看小结的符号表 + 定理速查表 + 选型表） & 25 分钟 & 写代码、选方法时回来查 \\
\bottomrule
\end{tabular}
\end{table}

\begin{note}
\textbf{本章记号与约定.}\;
本章工作在规划 / 控制的自然记号下，与 \cref{ch:planning,ch:control} 一致：状态 $\mathbf{x}\in\mathbb{R}^{n_x}$、控制 $\mathbf{u}\in\mathbb{R}^{n_u}$；离散动力学 $\mathbf{x}_{k+1}=f(\mathbf{x}_k,\mathbf{u}_k)$ 或线性 $\mathbf{x}_{k+1}=A\mathbf{x}_k+B\mathbf{u}_k$；状态约束集 $\mathcal{X}$、控制约束集 $\mathcal{U}$；时域长度 $N$。
凡涉及刚体姿态者，旋转 $\mathbf{R}\in\SO(3)$、位姿 $\mathbf{T}\in\SE(3)$，沿用 \cref{ch:lie} 的 $\SO,\SE,\so,\se$ 与右扰动约定（本章正文以平面 / 点质量 / 双积分器为主，姿态记号仅在引用足式控制时出现）。
\textbf{字母多义声明}（务必先读）：本章三个 $\alpha$ 含义不同——CBF 条件里的 $\alpha(\cdot)$ 是\emph{扩展类 $\mathcal{K}$ 函数}（\cref{sec:pr-cbf}）、Homothetic Tube 的 $\alpha_k$ 是\emph{管道缩放标量序列}（\cref{sec:pr-tube-gen}）、CVaR 的 $\alpha$ 是\emph{风险置信水平}（\cref{sec:pr-marc}）；$\Omega$ 在本章指 Tube MPC 的 RPI 集，与 POMDP 的观测空间无关；场景 / 情景记 $\omega_k$，意图概率 $P(\omega_k)$。所有多义字母均按所在小节的语境唯一确定。
\end{note}

% ════════════════════════════════════════════════════════════════
\section{不确定性从哪来：四类来源与两副面孔}
\label{sec:pr-uncertainty}

\paragraph{动机：MPC 的“完美未来”假设。}
\cref{sec:ctrl-mpc} 的名义 MPC 求解
\begin{equation}
  \min_{\mathbf{u}_{0:N-1}}\ \sum_{k=0}^{N-1}\ell(\mathbf{x}_k,\mathbf{u}_k)+\ell_f(\mathbf{x}_N)\quad\text{s.t.}\quad \mathbf{x}_{k+1}=A\mathbf{x}_k+B\mathbf{u}_k,\ \mathbf{x}_k\in\mathcal{X},\ \mathbf{u}_k\in\mathcal{U}.
  \label{eq:pr-nominal-mpc}
\end{equation}
它默认动力学\emph{精确}、未来\emph{唯一确定}。可现实里没有一处是确定的：模型不准、有扰动、他人意图不明。本节先把“不确定性”这个笼统词拆开——只有先认清\emph{是哪一类}不确定性，才能选对工具；否则用处理连续扰动的方法去对付离散意图（或反之），必然事倍功半。

\paragraph{反面：把所有不确定性当成一回事。}
一个常见错误是把“不确定性”一概当成“噪声方差”，统一用“加一圈安全裕度”应对。但“路口对向来车让还是不让”不是“让行程度的连续抖动”——它是两种\emph{定性不同}的未来。对它加裕度（按最坏情形避让）会导致 frozen；只赌最可能模态又会因预测翻转而决策振荡。混淆不确定性的\emph{类型}，是规控系统失败的常见根因。

\subsection{四类来源：认知 / 偶然 / 模型 / 环境}
\label{sec:pr-unc-types}

按来源，规控中的不确定性可分为四类（分类沿用不确定性规控的总论框架）：

\begin{description}
  \item[认知不确定性（epistemic）] 来自\emph{知识 / 数据不足}，原则上\emph{可消减}——多采数据、多观测就能降低。例：某地形的摩擦系数未知，但可通过试探估计。
  \item[偶然不确定性（aleatoric）] 来自系统\emph{固有随机性}，是消不掉的“噪声地板”——再多数据也降不下去。例：传感器的本底噪声、执行器的随机抖动。
  \item[模型不确定性（model）] 名义模型与真实系统的\emph{系统性偏差}，常\emph{可学习}。例：四旋翼高速飞行的气动残差、赛车的轮胎-路面交互——名义刚体模型解释不了这部分。
  \item[环境不确定性（environmental）] 来自\emph{外部世界的不可预测}，尤以\emph{他人意图}为代表——离散、多模态。例：他车让 / 不让、行人左转 / 右转。
\end{description}

\begin{insight}
四类来源在“能不能消”“怎么消”上分得很清，这直接决定方法选择：认知不确定性 → 主动感知 / 数据驱动学习（GP-MPC 的方差就是它的信号，\cref{sec:pr-gpmpc}）；偶然不确定性 → 概率 / 机会约束容忍（无法消除，只能在约束里留概率余量）；模型不确定性 → 学残差修正（GP / NN）；环境不确定性 → 分支 / 博弈规划（\cref{sec:pr-branch-form}）。\textbf{把“方差大”一律归因为“数据不够”是典型误区}——若方差主要来自偶然不确定性，再多采数据也没用（\cref{sec:pr-gpmpc} 会再遇到这个坑）。
\end{insight}

\subsection{两副面孔：连续小扰动 vs 离散多模态}
\label{sec:pr-unc-faces}

从\emph{处理方式}看，四类来源收敛成两副面孔，正对应本章两条主线：

\begin{itemize}
  \item \textbf{连续有界 / 概率扰动}（认知 + 偶然 + 模型）：围绕名义轨迹的“小邻域抖动”，是\emph{单峰}分布的方差。物理图像是“风把车吹偏一点点”。面对它，留一圈安全裕度就够——\emph{一个动作应对所有抖动}。这是\textbf{鲁棒规划与安全滤波}（\cref{sec:pr-tube,sec:pr-reach,sec:pr-gpmpc,sec:pr-cbf}）的领地。
  \item \textbf{离散多模态意图}（环境）：几条分开的、定性不同的未来，是\emph{多峰}（混合）分布的“哪个峰”。物理图像是“岔路口往左还是往右”。面对它，单一动作往往\emph{无法同时对所有峰最优}（抢行就该刹车、让行就该通过），你得“先走一个对两者都不坏的动作，等看清了再分别应对”。这是\textbf{分支 / 场景规划}（\cref{sec:pr-branch-form,sec:pr-mpdm,sec:pr-eudm,sec:pr-epsilon,sec:pr-marc}）的领地。
\end{itemize}

\begin{insight}
\textbf{两副面孔正交、可叠加}。一个真实安全系统常常是“分支处理离散意图 + 每条分支内用机会约束处理连续扰动 + CBF 兜瞬时安全底线”三者并用——意图是离散多峰，扰动是连续单峰，各管各的维度。\cref{sec:pr-branch-extra} 会把“分支 $\times$ 机会约束”的嵌套形式写出来，正是这一点的精确落地。
\end{insight}

\begin{pitfall}[概念误区：把离散意图塞进连续方差]
新手常想把“让 / 不让”也当成连续分布的一个大方差去处理。这会让单峰高斯假设崩掉、约束收紧失效——你无法用“一个均值 + 一个方差”同时表达“抢行（刹车）”和“让行（通过）”这两个\emph{对立}的最优动作。正确做法是\textbf{分层}：离散意图用分支（多峰各留一条独立策略），每峰内的连续扰动才用方差收紧（单峰）。这与高斯混合下“把多峰拆成多个高斯各自处理”同理。
\end{pitfall}

\begin{practice}
\begin{enumerate}
  \item 给四个场景分类（认知 / 偶然 / 模型 / 环境）并说明该用本章哪条主线：(a) 未知地面是冰还是干；(b) IMU 的随机游走噪声；(c) 机械臂关节摩擦未建模；(d) 前车直行还是右转下匝道。
  \item 解释为什么“对悲观预测加最大裕度”会导致 frozen，而“只赌最可能模态”会导致决策振荡——它们分别错在把离散意图当成了什么？
\end{enumerate}
\end{practice}

% ════════════════════════════════════════════════════════════════
\section{分支 / contingency 规划的基本形式化}
\label{sec:pr-branch-form}

\paragraph{动机：未来不是“一条”，而是“几种”。}
\cref{eq:pr-nominal-mpc} 的名义 MPC，骨子里假设未来是一条确定的轨迹。但他人意图的不确定性根本不是“小抖动”，而是\emph{离散的、分叉的}：对向来车要么让行、要么抢行；旁车道后车要么减速放你进、要么加速堵你。这些可能的未来是\emph{少数几个离散模态}，每个对应一条定性不同的演化；而且\emph{哪个模态会发生，往往要再过一会儿才能从观测看清}（车开始减速了，你才知道它要让你）。分支规划的全部出发点，就是把“赌最可能的单一未来”换成“为多个离散未来分别准备”。

\paragraph{反面：只赌最可能的未来，两个失败模式。}
设想最朴素的做法——用预测器给出每个他车的未来轨迹，当成确定障碍做一次普通单轨迹 MPC。两个典型失败：
\begin{itemize}
  \item \textbf{frozen robot（僵住）}：若预测“覆盖所有可能性”（他车可达集很大），单轨迹规划器发现整条路都被“可能占用”的区域堵死，找不到可行轨迹，机器人停下不动。预测越保守，越容易僵住。
  \item \textbf{决策振荡（policy oscillation）}：若预测“当前最可能”的单一模态，那么当观测把“最可能”从模态 A 翻到 B（他车细微动作让概率反转），最优解\emph{突变}——上一帧加速通过、这一帧猛刹让行。反复翻转即决策振荡，乘客体验极差且不安全。
\end{itemize}

\begin{insight}
\textbf{单轨迹规划的两难，根源是“过早承诺”}。frozen 是“对所有可能都承诺避让”（太保守、无解），振荡是“对当前最可能完全承诺”（太激进、易翻转）。两者都错在\emph{信息还不够时就把当前动作完全绑定到某个未来假设上}。分支规划的解法是\textbf{延迟承诺（commitment deferral）}：当前只走一段对所有候选未来都安全的“共享干”，把“到底按哪个未来走”推迟到分支点之后——那时观测多半已把真实意图揭示。这与“现金 / 期权”的实物期权（real option）、多阶段随机规划的 here-and-now + recourse 是同一原理，只是这里实现在轨迹空间里。
\end{insight}

\subsection{形式化：共享干 + 加权分支}
\label{sec:pr-branch-formal}

把直觉写成优化问题。在当前时刻 $t_0$，ego 面对 $K$ 个离散未来情景 $\{\omega_1,\dots,\omega_K\}$（如他车的 $K$ 种意图），每个有概率 $P(\omega_k)$（$\sum_k P(\omega_k)=1$，来自预测器或先验）。我们决策一棵\emph{轨迹树}：一段所有情景共用的“共享干” + 每个情景一条独立的“分支”\cite{hardy2013contingency}。

\begin{definition}[分支 / contingency 规划]\label{def:pr-contingency}
设共享干控制 $\mathbf{u}_{\mathrm{trunk}}$、各分支控制 $\{\mathbf{u}_k\}_{k=1}^K$，分支点 $t_{\mathrm{branch}}$。分支规划求解
\begin{equation}
\begin{aligned}
\min_{\mathbf{u}_{\mathrm{trunk}},\,\{\mathbf{u}_k\}}\quad & \sum_{k=1}^{K} P(\omega_k)\, J\big(\mathbf{u}_{\mathrm{trunk}},\,\mathbf{u}_k;\,\omega_k\big)\\
\text{s.t.}\quad
& \mathbf{u}(t_0:t_{\mathrm{branch}}) = \mathbf{u}_{\mathrm{trunk}} && (\text{所有 }\omega_k\text{ 相同——共享干 / 非预期约束})\\
& \mathbf{u}(t_{\mathrm{branch}}:T,\,\omega_k) = \mathbf{u}_k && (\text{每个 }\omega_k\text{ 独立优化——分支})\\
& \mathbf{x}_{k,t+1}=f(\mathbf{x}_{k,t},\mathbf{u}_{k,t},\omega_k),\ \ \mathbf{x}_{k,t}\in\mathcal{X}(\omega_k),\ \ \mathbf{u}_{k,t}\in\mathcal{U} && \forall k,t\\
& \mathbf{x}_{k,t_0}=\mathbf{x}_0 && \forall k.
\end{aligned}
\label{eq:pr-contingency}
\end{equation}
执行时只施加\emph{共享干的第一个控制} $\mathbf{u}_{\mathrm{trunk}}^{(0)}$，然后滚动重规划。
\end{definition}

逐项解读：
\begin{itemize}
  \item \textbf{场景加权代价} $\sum_k P(\omega_k) J(\cdot;\omega_k)$：优化所有未来的\emph{期望}代价，而非只优化某一个未来。
  \item \textbf{共享干约束}（核心）：$t_0$ 到 $t_{\mathrm{branch}}$ 之间的控制 $\mathbf{u}_{\mathrm{trunk}}$ 对所有情景\emph{必须相同}。这是\textbf{非预期约束（non-anticipativity constraint）}——在不确定性揭示前（$t<t_{\mathrm{branch}}$），你还分不清是哪个 $\omega_k$，所以此刻只能下一个对所有情景一致的动作。它保证短期决策一致（消除振荡）。
  \item \textbf{分支约束}：$t_{\mathrm{branch}}$ 之后，每个情景用各自独立的控制、在各自的动力学 / 约束 $f(\cdot,\omega_k),\mathcal{X}(\omega_k)$ 下优化，为每个未来留一条专门出路（化解 frozen）。
  \item \textbf{约束随情景变} $\mathcal{X}(\omega_k)$：障碍 / 可行域本身依赖情景——他车让行时占的空间和抢行时不同。
\end{itemize}

\begin{insight}
\textbf{非预期约束 = “信息不能预知未来”这条铁律的直接编码}。用多阶段随机规划的语言：决策在时间上展开，每步只知道“到目前为止观测到的信息”（信息过滤流 $\mathcal{F}_0\subseteq\mathcal{F}_1\subseteq\cdots$），决策 $\mathbf{u}_t$ 必须是 $\mathcal{F}_t$-可测的。在场景树上这翻译成：\emph{若两个场景在节点 $t$ 处尚未分开（共享同一历史），则它们在该节点的决策必须相同}。共享干就是“所有场景尚未分开的公共历史”，分支点 $t_{\mathrm{branch}}$ 是信息揭示（场景可区分）的时刻。这把分支规划牢牢锚定在随机规划的经典理论上——非预期约束不是工程 trick。
\end{insight}

\begin{pitfall}[概念误区：把分支规划当成“为每个场景各跑一个独立 MPC”]
新手想法：“不就是对 $K$ 个未来各算一条最优轨迹吗？分别跑 $K$ 次 MPC 就行。”后果：$K$ 条轨迹的“当前动作”互不相同，你不知道该执行哪一个；若按“最可能场景”那条执行，就退回了单轨迹规划的决策振荡。\textbf{根本原因}：漏掉了共享干（非预期约束）——分支规划的灵魂不是“分别优化”，而是“分支前必须共用一个当前动作”，正是这条耦合约束让“当前该做什么”有唯一、稳健的答案。\textbf{编程层面}的对应陷阱：非预期约束要加在\emph{控制}上（\verb|u[j][:, :N_branch] == u[0][:, :N_branch]|），不能只加在状态上——你执行的是控制不是状态，“此刻只能下一个动作”约束的对象必须是控制。检验：把 $t_{\mathrm{branch}}$ 设到 $0$（无共享干），看决策是否开始随预测翻转而振荡。
\end{pitfall}

\subsection{三个关键设计旋钮}
\label{sec:pr-branch-knobs}

\cref{eq:pr-contingency} 里有三个必须定的量，它们决定一个分支规划器的全部行为：

\begin{table}[H]\centering\small
\caption{分支规划的三个设计旋钮及各方法的取法}
\label{tab:pr-knobs}
\begin{tabular}{p{0.16\textwidth} p{0.30\textwidth} p{0.44\textwidth}}
\toprule
旋钮 & 含义 & 各方法怎么定 \\
\midrule
$t_{\mathrm{branch}}$ 分支点 & 共享干多长、何时分叉 & MPDM / EUDM：\textbf{固定}（EUDM 用 $8\,$s 时域、$2\,$s 一个策略节点、树深 4）；MARC：\textbf{动态}（按场景轨迹“开始显著分叉”处，用场景级偏差函数度量） \\
$K$ 分支数 & 枚举几个未来 & MPDM：$O(|\Pi_{\mathrm{ego}}|\cdot|\Pi_{\mathrm{other}}|^{n})$（组合，易爆炸）；EUDM：CFB 裁到 $3$--$8$ 个 \\
$P(\omega_k)$ 场景概率 & 各未来的权重 & 来自预测器输出 / 领域先验 / 对他车意图的贝叶斯更新 \\
\bottomrule
\end{tabular}
\end{table}

\textbf{为什么 $t_{\mathrm{branch}}$ 最关键}：它太早（共享干太短）$\to$ 退化成 $K$ 个几乎独立的 MPC，短期一致性弱、易振荡；太晚（共享干太长）$\to$ 所有情景被迫共用很长一段动作，过度保守（回到 frozen 边缘）。“在信息真正揭示的时刻分叉”才最优——这正是 MARC 相对 MPDM/EUDM 的核心改进（\cref{sec:pr-marc}）。

\subsection{问题规模走查：分支比单轨迹大多少}
\label{sec:pr-branch-size}

把“规模 $\sim K\times$（但共享干减少增量）”用具体数字坐实。设时域 $N=30$ 步、状态 $n_x=2$（位置 $s$、速度 $v$）、控制 $n_u=1$（加速度 $a$）、分支数 $K=2$、共享干 $N_{\mathrm{trunk}}=10$。
\begin{itemize}
  \item \textbf{单轨迹 MPC}：控制变量 $N\cdot n_u=30$ 个。
  \item \textbf{分支规划}：控制变量 $=\underbrace{N_{\mathrm{trunk}}}_{\text{共享干}}+\underbrace{K\cdot(N-N_{\mathrm{trunk}})}_{\text{各分支}}=10+2\times20=50$ 个（\textbf{不是} $K\cdot N=60$——共享干那 10 步只算一次）。
\end{itemize}
对比：控制变量 $50$ vs 单轨迹 $30$，是 $1.67\times$（而非 $K=2$ 倍）——差的 $0.33\times$ 正是共享干省下的。共享干越长增量越小：$N_{\mathrm{trunk}}=20$ 时控制变量 $=20+2\times10=40$（仅 $1.33\times$）；$N_{\mathrm{trunk}}=0$ 时 $=2\times30=60$（满 $2\times$，退化成独立两段）。

\begin{insight}
\textbf{分支规划的成本几乎全在“分支数 $K$”，收益几乎全在“共享干”}。规模随 $K$ 近似线性增长，所以控制 $K$（EUDM 的 CFB 干这个）是工程可行性的关键；而真正解决 frozen / 振荡两难的，是共享干那条非预期约束。记住：\emph{分支是为了不漏掉任何未来，共享干是为了不过早承诺任何未来}。后续所有方法（MPDM / EUDM / EPSILON / MARC / multi-stage NMPC）都只是在回答“$K$ 个情景怎么选、怎么压、分支点放哪、代价怎么加权”——骨架就是 \cref{eq:pr-contingency} 这个“共享干 + 加权分支”。
\end{insight}

\begin{practice}
\begin{enumerate}
  \item \textbf{[形式化]} 把 \cref{eq:pr-contingency} 退化到 $K=1$（只有一个未来），证明它就是普通单轨迹 MPC（\cref{eq:pr-nominal-mpc}）。再退化到 $t_{\mathrm{branch}}=T$（共享干贯穿全程），说明此时 $K$ 个场景被迫共用整条轨迹——这对应对场景集做某种“期望 / 最坏”处理。
  \item \textbf{[规模]} 设 $N=40$、$K=3$、$N_{\mathrm{trunk}}=15$、$n_u=2$。算分支规划的控制变量数，与单轨迹和满 $K\times$ 各比多少倍？说明共享干长度如何调节这个倍数。
  \item \textbf{[思考题]} 用“承诺时机”一个概念统一解释 frozen robot 与决策振荡，并说明共享干 + 分支如何同时治好这两个病。
\end{enumerate}
\end{practice}

% ════════════════════════════════════════════════════════════════
\section{MPDM：多策略决策}
\label{sec:pr-mpdm}

\paragraph{动机：他车的未来不是任意轨迹，而是少数几种“语义策略”。}
\cref{eq:pr-contingency} 里 $K$ 个情景 $\omega_k$ 是抽象的“离散未来”。\textbf{MPDM（Multi-Policy Decision Making，多策略决策）}\cite{cunningham2015mpdm} 给了它一个极有用的具体化：把每辆车（含 ego 和他车）的未来建模为\emph{从一个小的、语义明确的策略集合里选一个}。人开车不是在连续控制空间任意游走，而是在执行少数几个语义策略：跟车、左变道、右变道、加速、减速。MPDM 的洞察是——既然如此，\emph{与其预测他车的连续轨迹，不如枚举它在执行哪个语义策略}。这把“无穷维的轨迹预测”降成“有限个离散策略的概率分配”，正好对应 $K$ 个情景。

\subsection{算法：策略枚举 + 闭环前向仿真}
\label{sec:pr-mpdm-algo}

MPDM 的核心循环：枚举 ego 的候选策略，对每个 ego 策略、在他车各候选策略下做\emph{闭环前向仿真}，按他车策略概率加权得到该 ego 策略的期望代价，最后选代价最小的 ego 策略\cite{cunningham2015mpdm}。

\begin{algo}[MPDM 策略枚举与评估]\label{alg:pr-mpdm}
对 ego 策略集 $\Pi_{\mathrm{ego}}$ 与他车策略集 $\Pi_{\mathrm{other}}$（及其概率 $P(\pi_{\mathrm{other}})$）做双重循环：每个 ego 策略对每个他车策略闭环仿真一段、算代价，按他车策略概率加权，取期望代价最小者。
\end{algo}

\begin{codebox}[text]{MPDM 核心循环}
for pi_ego in Pi_ego:
    for pi_other in Pi_other:          # or only top-k key vehicles
        # closed-loop forward sim 2-5s: ego runs pi_ego, others run pi_other,
        # both react via IDM/MOBIL (interaction is simulated, not replayed)
        rollout = closed_loop_sim(pi_ego, pi_other, horizon=T_sim)
        score[pi_ego][pi_other] = cost(rollout)   # safety + efficiency + comfort
    J[pi_ego] = sum_over_pi_other( P(pi_other) * score[pi_ego][pi_other] )
pi_star = argmin_{pi_ego} J[pi_ego]
execute first segment of pi_star, then receding-horizon replan
\end{codebox}

这正是 \cref{eq:pr-contingency} 的一个具体实例：$\omega_k=\pi_{\mathrm{other}}$，$P(\omega_k)=P(\pi_{\mathrm{other}})$，代价 $J$ 来自闭环前向仿真。MPDM 的“共享”体现在它\emph{选一个语义策略}（而非一条轨迹）——这个策略本身在短期内对各种他车反应都给出一致的动作。

\subsection{闭环前向仿真的两个反应式模型}
\label{sec:pr-mpdm-models}

MPDM 的前向仿真要让车“按策略自动驾驶”，靠两个经典微观交通模型。

\textbf{IDM（Intelligent Driver Model，智能驾驶员模型）}\cite{treiber2000idm} 管纵向（跟车），给出加速度作为“自己想多快”和“前车有多近”的函数：
\begin{equation}
  \dot v = a\left[1-\left(\frac{v}{v_0}\right)^{4}-\left(\frac{s^*(v,\Delta v)}{s}\right)^{2}\right],
  \qquad
  s^*(v,\Delta v)=s_0+\max\!\Big(0,\ vT+\frac{v\,\Delta v}{2\sqrt{ab}}\Big),
  \label{eq:pr-idm}
\end{equation}
五参数：期望速度 $v_0$、最大加速度 $a$、舒适减速度 $b$、最小停车间距 $s_0$、期望车头时距 $T$；$s$ 是与前车实际间距、$\Delta v=v-v_{\mathrm{lead}}$ 是接近速度。直觉：括号里第一项 $(v/v_0)^4$ 是“自由驾驶”项（没车挡时趋向 $v_0$），第二项 $(s^*/s)^2$ 是“交互”项（前车越近、相对速度越大，$s^*$ 越大、减速越狠）；$s\to s^*$ 时两项平衡、加速度趋零（稳定跟车）。一条光滑函数自然涵盖四个驾驶 regime——自由流（$s\gg s^*$，趋向 $v_0$）、接近（$\Delta v>0$，相对速度项使 $s^*$ 增大，\emph{预判刹车}）、跟随（$s\approx s^*$，稳定）、紧急（$s\ll s^*$，急刹）。

\textbf{MOBIL（Minimizing Overall Braking Induced by Lane changes，换道模型）}\cite{kesting2007mobil} 管横向，用 IDM 算“换道前后各车的加速度变化”，换道当且仅当同时满足：
\begin{align}
  \text{安全准则}\;&: & a^{\mathrm{new}}_{\mathrm{back}}&\ge -b_{\mathrm{safe}}, \label{eq:pr-mobil-safe}\\
  \text{激励准则}\;&: & \underbrace{(a^{\mathrm{new}}_{\mathrm{ego}}-a^{\mathrm{old}}_{\mathrm{ego}})}_{\text{自己的收益}}+p\underbrace{\big[(a^{\mathrm{new}}_{\mathrm{others}}-a^{\mathrm{old}}_{\mathrm{others}})\big]}_{\text{对相关邻车的影响}}&>\Delta a_{\mathrm{th}}, \label{eq:pr-mobil-incentive}
\end{align}
其中 $p\in[0,1]$ 是\emph{礼让系数（politeness）}（$p=0$ 自私、$p=1$ 完全考虑邻车），$\Delta a_{\mathrm{th}}$ 是换道阈值（避免频繁变道），$b_{\mathrm{safe}}$ 是目标车道后车可承受的最大减速度。

\begin{insight}
\textbf{MPDM 的前向仿真用“行为模型”代替“轨迹预测”}。IDM / MOBIL 不预测某车的具体轨迹，而是给一个\emph{策略}（如“跟车”）配一套反应式动力学，让车在仿真里自己跑出轨迹。两个好处：(1) \textbf{交互闭环}——ego 和他车在仿真里相互反应，自动体现“ego 减速则他车跟进”；(2) \textbf{语义可解释}——每条仿真轨迹对应明确的策略组合，便于调试审计。代价是 IDM / MOBIL 是简化模型，真实司机未必严格遵守（EPSILON 后来加“安全机制”兜底这种模型失配，\cref{sec:pr-epsilon}）。
\end{insight}

\begin{insight}
\textbf{MOBIL 的“安全准则一票否决” vs “激励准则加权权衡”}。\cref{eq:pr-mobil-safe} 是\emph{硬门槛}（不满足直接不换，类似 CBF 硬安全约束），\cref{eq:pr-mobil-incentive} 是\emph{软权衡}（收益超阈值才换，类似代价里的偏好项）。这种“硬安全门槛 + 软收益权衡”二层结构在机器人决策里反复出现（CBF-QP 的硬约束 + 软目标、机会约束的硬违约上限 + 软效率代价）——安全永远优先于效率。
\end{insight}

\subsection{为什么闭环对 merge 场景是决定性的}
\label{sec:pr-mpdm-merge}

并线（merge）是分支规划文献的标志性场景，也是“开环必败、闭环必需”最清楚的例子\cite{cunningham2015mpdm}。\textbf{场景}：ego 在汇入车道想并入主路，主路目标间隙后方有车 B；ego 有 \{加速并入, 减速等待\}，B 有 \{让行, 不让\}。

\textbf{开环仿真为什么必败}：开环把 B 当“按预测轨迹走的录像”——B 的轨迹固定、不随 ego 反应。于是若预测 B“匀速不让”，ego 算出“间隙不够”→ 减速等待 → 但 B 其实会让时，ego 白白错过机会（过保守 / frozen 边缘）；开环永远看不到“ego 一压过去、B 就松油门让行”这种\emph{交互涌现}的间隙。

\textbf{闭环仿真为什么必需}：闭环让 ego 和 B 都按各自策略用 IDM 互相反应——ego“加速并入”+ B“让行”意图：B 检测到前方切入、自动减速让出空间 → 并入成功、代价低；ego“加速并入”+ B“不让”意图：B 不减速甚至加速 → 间隙不够、安全罚高。于是 MPDM 在闭环下能正确区分两个分支，按概率加权选策略。

\begin{insight}
\textbf{merge 的本质是“间隙是协商出来的，不是观测到的”}。开环假设间隙外生固定（B 的录像），ego 只能被动接受；闭环承认间隙是 ego 与 B 交互\emph{协商}的产物（ego 的动作改变 B 的反应）。分支规划 + 闭环仿真因此能表达“试探性前压 → 观察 B 反应 → 按反应分支”——这正是人类司机并线时做的事。
\end{insight}

\subsection{复杂度：MPDM 为什么不可扩展}
\label{sec:pr-mpdm-complexity}

MPDM 的双重循环里，外层是 ego 策略（$|\Pi_{\mathrm{ego}}|$ 个），内层枚举他车策略组合。若场景有 $n$ 辆相关他车、每辆 $|\Pi_{\mathrm{other}}|$ 个策略，他车联合策略组合数是 $|\Pi_{\mathrm{other}}|^{n}$，于是总仿真次数
\begin{equation}
  O\big(|\Pi_{\mathrm{ego}}|\cdot|\Pi_{\mathrm{other}}|^{n}\big).
  \label{eq:pr-mpdm-complexity}
\end{equation}
举例：$|\Pi_{\mathrm{ego}}|=5,|\Pi_{\mathrm{other}}|=3,n=3$ → $5\times3^3=135$ 次；$n=5$ → $5\times3^5=1215$ 次。\textbf{指数项 $|\Pi_{\mathrm{other}}|^{n}$ 是杀手}——他车一多就炸。这正是 EUDM（\cref{sec:pr-eudm}）要解决的问题。

\begin{pitfall}[陷阱：以为 MPDM 预测精确轨迹 / 用开环仿真]
两个常见误区。其一：以为 MPDM 内部有个高精度轨迹预测器——其实 MPDM 预测的是\emph{离散策略的概率} $P(\pi_{\mathrm{other}})$（$\sum_\pi P(\pi)=1$），轨迹是用 IDM / MOBIL 在前向仿真里“跑”出来的。把预测器输出对接成“他车各策略的概率”，而非位置序列。其二：用开环（非交互）仿真评估策略——会丢掉交互，等于假设他车对 ego 毫无反应；在并线场景里开环会 frozen（挤不进），闭环才能在后车让行时并入。另外，实现 IDM 时 $s^*$ 必须取 $\max(0,\cdot)$ 且保留相对速度项 $\frac{v\Delta v}{2\sqrt{ab}}$，否则跟车迟钝、易追尾或出现负间距穿模。
\end{pitfall}

\begin{practice}
\begin{enumerate}
  \item \textbf{[实现 IDM]} 实现 IDM 单车道跟车：前车按给定速度曲线（含一次急刹）行驶，本车用 \cref{eq:pr-idm} 跟随。画速度 / 间距曲线，验证间距始终 $\ge s_0$、急刹时本车平滑减速。扫 $T$（车头时距）观察跟车松紧。
  \item \textbf{[实现 MPDM]} 在简化并线场景里实现最小 MPDM（\cref{alg:pr-mpdm}）：ego 有 \{加速并入, 减速等待\}，主道后车有 \{让行, 不让\}，用 IDM 做 $3\,$s 闭环前向仿真，按后车策略概率加权选 ego 策略。验证后车大概率让行时 ego 选并入、大概率不让时选等待。
  \item \textbf{[复杂度]} 给定 $|\Pi_{\mathrm{ego}}|=5,|\Pi_{\mathrm{other}}|=4$，分别算 $n=2,4,6$ 辆他车时的仿真次数（\cref{eq:pr-mpdm-complexity}）。若每次仿真 $1\,$ms，哪个 $n$ 开始超过 $100\,$ms 实时预算？
\end{enumerate}
\end{practice}

% ════════════════════════════════════════════════════════════════
\section{EUDM：引导分支把指数压回实时}
\label{sec:pr-eudm}

\paragraph{动机：把 MPDM 的指数爆炸压回实时。}
\cref{eq:pr-mpdm-complexity} 的 $O(|\Pi_{\mathrm{ego}}|\cdot|\Pi_{\mathrm{other}}|^n)$ 随他车数 $n$ 指数爆炸。\textbf{EUDM（Efficient Uncertainty-aware Decision-making，高效不确定性感知决策）}\cite{zhang2020eudm} 用两个机制把它压回实时，同时保持策略枚举思想。它把决策形式化为一个 POMDP（部分可观马尔可夫决策过程，他车意图是隐变量），但\emph{不直接求解} POMDP（那会爆炸），而是用领域知识“引导分支”。\pz{POMDP 的形式化与点基求解（PBVI/SARSOP）属决策理论范畴，本章只用其语言、不展开求解；如本书后续设 POMDP 专章可在此 \texttt{cref} 桥接。}

两个朴素路线各自怎么崩，就懂了“引导”的必要：\emph{精确求解 POMDP} 的复杂度随状态 / 观测 / 时域指数爆炸，自驾的连续状态 + 长时域远超精确求解器能力；\emph{MPDM 式全组合枚举} 则如 \cref{eq:pr-mpdm-complexity} 所示随车数指数爆炸。EUDM 在 action space（ego 动作）和 intention space（他车意图）\emph{两个空间都用领域先验引导分支}——只展开“值得展开”的少数分支。

\subsection{机制一：DCP-Tree（领域专属闭环策略树）}
\label{sec:pr-eudm-dcp}

MPDM 在\emph{单个时刻}对 ego 选一个策略，但真实驾驶是\emph{策略序列}——“先跟车 $1.5\,$s，再左变道”。若每个时间节点都允许 ego 切换到任意策略，策略序列数仍指数爆炸（$|\Pi_{\mathrm{ego}}|^{\text{深度}}$）。DCP-Tree（Domain-specific Closed-loop Policy Tree）用一条领域先验大幅剪枝：\textbf{在一个规划周期内，ego 最多切换一次策略}\cite{zhang2020eudm}。理由很实在——若真需要复杂的策略变化（“跟车→变道→再跟车”），可靠\emph{高频重规划}分多个周期实现，单周期内不必考虑多次切换。EUDM 用 $8\,$s 规划时域、每 $2\,$s 一个策略更新节点（树深 4），把单周期 ego 策略序列数从 $|\Pi_{\mathrm{ego}}|^4$ 压到线性级。

\begin{insight}
\textbf{DCP-Tree 用“一周期一次切换”把指数压成线性}。关键是“决策频率”与“策略复杂度”可解耦——你不需在一个周期里规划全部复杂机动，因为滚动重规划（每几十毫秒一次）会不断给你新机会切换。这是典型的“用时间换空间”：把“一次想清楚所有切换”换成“每周期只想一次切换 + 高频重规划”，与 MPC 滚动时域一脉相承。
\end{insight}

\subsection{机制二：CFB（条件聚焦分支）}
\label{sec:pr-eudm-cfb}

DCP-Tree 解决了 ego 侧，他车侧呢？CFB（Conditional Focused Branching）的洞察：\emph{绝大多数他车与 ego 的当前决策无关，不值得为它们分支}\cite{zhang2020eudm}。CFB 对每条 ego 策略序列，只在“危险交互”处分支他车意图，其余他车一律用其\emph{最可能}策略（不分支）。“危险交互”指：在这条 ego 策略下，某个他车的不同意图（让 / 不让）会显著改变 ego 的安全性或代价。CFB 通过一个初步的（用最可能意图的）前向仿真识别出“关键交互车”，只对这少数几辆分支。于是分支数从“所有他车组合”$|\Pi_{\mathrm{other}}|^n$ 压到“少数关键交互车组合”$|\Pi_{\mathrm{other}}|^{n_{\mathrm{key}}}$（$n_{\mathrm{key}}=1\sim2$），通常 $3$--$8$ 个分支。

\subsection{联合复杂度：DCP-Tree × CFB}
\label{sec:pr-eudm-joint}

两机制相乘得到 EUDM 的最终复杂度\cite{zhang2020eudm}：
\begin{equation}
  O\big(\underbrace{d\cdot|\Pi_{\mathrm{ego}}|}_{\text{DCP-Tree（线性）}}\times\underbrace{|\Pi_{\mathrm{other}}|^{n_{\mathrm{key}}}}_{\text{CFB（指数的指数被砍）}}\big),
  \label{eq:pr-eudm-complexity}
\end{equation}
对比 MPDM 的 $O(|\Pi_{\mathrm{ego}}|\cdot|\Pi_{\mathrm{other}}|^n)$：ego 侧从（潜在的）$|\Pi_{\mathrm{ego}}|^d$ 指数降到 $d\cdot|\Pi_{\mathrm{ego}}|$ 线性（DCP-Tree 省），他车侧把指数的指数从 $n$ 砍到 $n_{\mathrm{key}}$（CFB 省）。数量级体感：$n=5,|\Pi_{\mathrm{other}}|=3$ 时 MPDM 的他车组合是 $3^5=243$，CFB 取 $n_{\mathrm{key}}=2$ 降到 $3^2=9$——\textbf{27 倍压缩}，且压缩比随 $n$ 增大急剧拉大。正是这两刀，让分支规划从“几辆车就超时”变成“$50$--$100\,$ms 内可解”。

\begin{insight}
\textbf{DCP-Tree 砍“底数的幂”、CFB 砍“指数本身”}。两机制压的是复杂度公式的不同部位——砍指数（CFB，把 $n\to n_{\mathrm{key}}$）比砍底数（DCP-Tree）收益更大，因为指数 $n$ 是爆炸的主因。记住这个“砍指数优先于砍底数”的直觉，你在任何组合优化里都知道该先攻哪里。CFB“只把计算花在关键交互处”的思想，与“风险预算分给紧约束”“多障碍 CBF 只激活最近几个”是同一种工程铁律。
\end{insight}

\begin{pitfall}[陷阱：把 DCP-Tree 与 CFB 的职责搞混 / 怕错过复杂机动]
其一：以为 DCP-Tree 在枚举他车策略——错。DCP-Tree 展开的是\emph{ego 自己的策略序列}（纵向 = 时间轴上的策略链），他车意图分支是 \emph{CFB} 干的（横向 = 在危险点）。画出树：纵向是 ego 策略序列、横向才是他车意图分支。其二：以为“一周期一次切换”会让 EUDM 错过“跟车→变道→回正”这种多次切换的机动——忽略了\emph{滚动重规划}：单周期一次切换 + 每几十毫秒重规划，多个周期串起来可实现任意复杂序列，表达能力没丢，只是分摊到时间轴上。
\end{pitfall}

\begin{practice}
\begin{enumerate}
  \item \textbf{[DCP-Tree]} 给定 ego 策略集 \{跟车, 左变道, 右变道\}、树深 4、约束“一周期 $\le 1$ 次切换”。枚举所有合法 ego 策略序列、数有多少条；对比“允许任意切换”的 $3^4=81$ 条，DCP-Tree 剪掉了多少？
  \item \textbf{[复杂度对比]} 用 \cref{eq:pr-mpdm-complexity,eq:pr-eudm-complexity}，对 $n=2,4,6,8$ 辆车（$|\Pi_{\mathrm{other}}|=3,n_{\mathrm{key}}=2$）画 MPDM 与 EUDM 的分支数曲线，定量说明压缩比随 $n$ 的变化。
  \item \textbf{[思考题]} CFB 靠“最可能意图的初步仿真”识别关键车。若某他车真实意图恰好是小概率的危险意图（初步仿真没把它当关键车），会发生什么？这暴露了 CFB 的什么假设？（联系 \cref{sec:pr-marc} 为何要改进分支策略）
\end{enumerate}
\end{practice}

% ════════════════════════════════════════════════════════════════
\section{EPSILON：从行为决策到完整分层系统}
\label{sec:pr-epsilon}

\paragraph{动机：EUDM 只给“策略”，谁来给“轨迹”？}
EUDM 是\emph{行为决策层}——输出“最优 ego 策略序列 + 粗轨迹参考”，但不给可执行的光滑轨迹。一个能上车的系统还需\emph{运动规划层}把策略落成满足动力学、避障、舒适的轨迹。\textbf{EPSILON（Efficient Planning System for Automated Vehicles in Highly Interactive Environments）}\cite{ding2021epsilon} 把 EUDM（行为层）与 \textbf{SSC}（运动层）拼成完整、开源、实车验证的分层系统。

\subsection{架构：行为决策层 + 运动规划层}
\label{sec:pr-epsilon-arch}

\begin{center}
\small
\begin{tikzpicture}[
  node distance=4mm, >=Stealth,
  box/.style={draw, rounded corners, align=center, font=\small, inner sep=4pt, fill=main!6, draw=main!60!black, text width=0.62\textwidth},
  ln/.style={->, thick, gray!60!black}]
\node[box, fill=second!6, draw=second!60!black] (env) {环境理解（感知 / 预测 / 定位）\\ {\scriptsize 输出：他车状态 + 意图先验 $P(\omega_k)$}};
\node[box, below=6mm of env] (eudm) {\textbf{行为决策层（EUDM）}\\ POMDP 形式化 + 引导分支（DCP-Tree + CFB）\\ {\scriptsize 输出：最优 ego 策略序列 + 粗轨迹参考}};
\node[box, below=6mm of eudm, fill=third!8, draw=third!60!black] (ssc) {\textbf{运动规划层（SSC）}\\ 构造时空语义走廊（凸多面体序列）→ 走廊内 QP 求光滑轨迹\\ {\scriptsize 输出：光滑可执行轨迹}};
\node[below=5mm of ssc, font=\small] (ctrl) {控制 / 执行};
\draw[ln] (env)--(eudm); \draw[ln] (eudm)--(ssc); \draw[ln] (ssc)--(ctrl);
\end{tikzpicture}
\end{center}

\textbf{为什么要分层}：行为决策（离散、组合、博弈味重）和运动规划（连续、光滑、几何约束）是两类不同问题，强行塞进一个优化里又非凸又难解。分层让行为层用枚举式策略树处理“离散意图”，运动层用凸优化处理“连续轨迹”——各用最适合的工具。这与 CBF 兜底 MPC（\cref{sec:pr-cbf-mpc}）、足式栈“MPC 给质心轨迹、WBC 给关节力矩”（\cref{sec:ctrl-legged}）是同一种分层哲学。

\subsection{SSC：时空语义走廊}
\label{sec:pr-epsilon-ssc}

SSC（Spatio-temporal Semantic Corridor）\cite{ding2019ssc} 是 EPSILON 的运动规划层。在 EUDM 选定的策略下（如“在他车后方并入”），SSC 把“该策略允许 ego 走的时空区域”编码成一串\emph{凸多面体}（在位置-时间空间里），即“语义走廊”，然后在走廊内解一个 \textbf{QP} 求光滑轨迹（位置 / 速度 / 加速度连续、在走廊内、满足动力学）。概念上分三步：(1) 把 ego 可走区域按选定策略与各障碍的时空占据切成一串\emph{语义立方体}（每个是一段时间内允许占据的位置区间，明确“在哪辆车前 / 后、哪条车道”）；(2) 把相邻立方体连成\emph{凸多面体序列}（走廊）；(3) 走廊内最小化光滑性（jerk / 加速度）+ 跟踪代价的 QP。因约束（凸多面体）与代价（二次）都凸，QP 快而可靠。

\begin{insight}
\textbf{EPSILON = “离散决策定通道 + 连续优化填轨迹”}。难的非凸部分（该走哪条通道、避哪辆车的哪一侧）交给离散的策略枚举；凸的部分（通道内怎么走最光滑）交给 QP。这解释了为什么\emph{纯连续优化做交互式驾驶会陷入局部最优}：“从障碍左侧绕”和“从右侧绕”是\textbf{同伦类（homotopy class）}不同、拓扑上不连续的两条轨迹——你无法在不穿过障碍的前提下把一条连续形变成另一条。纯连续轨迹优化（一个 NLP / QP 做梯度下降）只能在\emph{当前同伦类内}找局部最优，跨不过障碍这道墙。分支 / 分层显式枚举不同的离散选择（每个对应一个同伦类），再在选定的\emph{单一}同伦类内做凸优化——分层不是可选技巧，而是\emph{应对离散拓扑选择的必需结构}。
\end{insight}

\subsection{ego-centric belief 与安全机制}
\label{sec:pr-epsilon-belief}

EPSILON 还加了两个工程优化\cite{ding2021epsilon}：\textbf{ego-centric belief update}——不对场景里所有车维护完整意图信念，只对 ego 当前决策真正关心的交互车维护 belief（与 CFB“只对关键交互车分支”一脉相承，信念维护也是有成本的）；\textbf{带安全机制的驾驶员模型}——当 IDM / MOBIL 这类先验模型与观测明显不符时，安全机制接管，对冲“真实他车未必严格按 IDM 跟车”的模型失配。EPSILON 开源于 \texttt{HKUST-Aerial-Robotics/EPSILON}，是少数把“行为决策 + 运动规划”完整闭环、且在真实城市车流验证过的开源交互式规划系统之一。\pz{EPSILON 的开源协议、实车验证的具体城市与场景规模，及“带安全机制的驾驶员模型”的具体形式，待联网核实原 T-RO 论文。}

\begin{pitfall}[陷阱：以为 EPSILON 是端到端单一优化 / 给所有车都维护信念才严谨]
其一：试图写一个“既选策略又出轨迹”的联合优化——又非凸又难解。EPSILON 的关键正是\emph{分层解耦}：行为层（离散枚举）输出语义策略，运动层（凸 QP）据此构造凸走廊再 QP。能清楚说出“哪一层处理离散、哪一层处理连续”才算理解。其二：以为对全场景所有车建信念才“严谨”——忽略了 ego-centric 洞察：绝大多数车与 ego 当前决策无关，为它们维护信念是浪费。严谨性应服务于“对 ego 决策有影响的不确定性”。
\end{pitfall}

\begin{practice}
\begin{enumerate}
  \item \textbf{[SSC 直觉]} 在“在前车后方并入”场景里，手工画出该策略对应的时空语义走廊（位置-时间空间的凸多面体序列），说明为什么在走廊内做 QP 是凸的。
  \item \textbf{[同伦类]} 给一个含单个动态障碍的二维场景，画出“绕左”与“绕右”两条轨迹，说明它们为何不同伦、以及纯梯度下降为何只能停在初值所在的那一类。
  \item \textbf{[思考题]} EPSILON 用“带安全机制的驾驶员模型”对冲 IDM / MOBIL 失配。这本质上在处理 \cref{sec:pr-unc-types} 的哪一类不确定性？它和鲁棒规划“对最坏扰动留余量”在思想上有何相通？
\end{enumerate}
\end{practice}

% ════════════════════════════════════════════════════════════════
\section{MARC：动态分支点 + 风险感知}
\label{sec:pr-marc}

\paragraph{动机：固定分支点的两个毛病。}
MPDM / EUDM / EPSILON 都在\emph{固定时刻}分支（如 EUDM 每 $2\,$s 一个策略节点）。\cref{sec:pr-branch-knobs} 已指出 $t_{\mathrm{branch}}$ 是最关键旋钮，固定它有两个毛病\cite{li2023marc}：(1) \textbf{分早了浪费 / 分晚了危险}——真实场景里“意图何时揭示”是变化的，有时一秒内表态（该早分支）、有时好几秒才看清（晚分支才合理）；固定分支点要么过早分叉（共享干太短、不够稳）要么过晚（共享干太长、过保守）。(2) \textbf{只看期望、不看尾部风险}——按场景概率加权\emph{期望}代价，会把“小概率但灾难性”的尾部（他车突然抢行导致碰撞）平均掉。\textbf{MARC（Multipolicy and Risk-aware Contingency Planning）}\cite{li2023marc} 针对这两点给出两项关键进化。

\subsection{进化一：场景级 divergence 动态决定分支点}
\label{sec:pr-marc-branch}

MARC 不在固定时刻分支，而是观察\emph{多个场景的轨迹何时开始显著分叉}，在那里动态设分支点。原文用一个\emph{场景级偏差函数} $\mathcal{M}(s_i,s_j,t)$ 度量任意一对场景 $(s_i,s_j)$ 在 $t$ 时刻 ego 状态的偏差，并把分支点取为\emph{所有场景对仍处于偏差阈值 $\theta$ 之内的最晚时刻}（在那之前各场景尚不可区分，故还不必分）：
\begin{equation}
  t_{\mathrm{branch}}=\max\Big\{t:\ \mathcal{M}(s_i,s_j,t)<\theta\ \ \forall (s_i,s_j)\in\mathcal{S}\times\mathcal{S}\Big\},
  \label{eq:pr-marc-tbranch}
\end{equation}
$\theta$ 是状态分叉阈值。\pz{原文（li2023marc）只给出 $\mathcal{M}$ 是“度量场景对间 ego 状态偏差”的函数、未固定其具体形式，并\emph{未使用} Wasserstein 距离；本书早先把该度量写成 Wasserstein 距离是不准确的，现按原文改回通用偏差函数。对轨迹分布若要量化分叉，Wasserstein 是一种可选度量，但非 MARC 原文所采用。}直觉：分支点应放在“信息真正揭示、各场景不得不分开”的时刻——两个场景前 $3\,$s 最优轨迹几乎重合（无论让不让 ego 前 $3\,$s 都该这么走），共享干就该到 $3\,$s；若 $1\,$s 就分开了，共享干只到 $1\,$s。\emph{让数据决定分支点}，而非先验固定。

对确定性轨迹，最简的偏差度量就是位置差 $\mathcal{M}(t)=|s_A(t)-s_B(t)|$。一个数值示意（意图 A 让行在 $1.5\,$s 后减速、意图 B 通过在 $1.5\,$s 后加速，之前最优动作相同、轨迹重合）：$t\in[0,1.5]$ 时 $\mathcal{M}=0$（共享无损区）；$t=2.0$ 时 $\mathcal{M}=0.875$、$t=2.5$ 时 $\mathcal{M}=2.625$、$t=3.0$ 时 $\mathcal{M}=5.250$。取阈值 $\theta=1.0$，最后一个仍满足 $\mathcal{M}<\theta$ 的时刻是 $t=2.0\,$s、下一时刻 $t=2.5\,$s 已越阈值——\emph{动态分支点落在二者之间（约 $2.5\,$s）}。对比 EUDM 固定在 $2.0\,$s 分（此时 $\mathcal{M}=0.875<\theta$，轨迹还没真正分开，分早了）。

\begin{insight}
\textbf{MARC 把分支点从“超参数”变成“由场景分叉自动涌现的量”}。理想分支点是“信息揭示时刻”，而它\emph{可观测}——就是“各场景最优轨迹开始显著分叉”的时刻（在那之前不同未来对应的最优当前动作一样，所以还不必分）。用场景级偏差函数 $\mathcal{M}$ 量化“分叉”，分支点就从手调超参变成自适应量。$\mathcal{M}(t)=0$ 的区间是“共享无损区”，延长共享干到此处末端不损失任何最优性——\emph{固定分支点是猜共享干该多长，动态分支点是测出来}。\pz{动态分支点示意里的具体数值（0.875/2.625/5.250 等）来自玩具算例、非论文实验数据，仅作教学示意；引用 MARC 时这些数字不应作为论文结论引用。}
\end{insight}

\subsection{进化二：CVaR 风险感知的双层优化}
\label{sec:pr-marc-cvar}

MARC 把“期望代价”换成 \textbf{CVaR（Conditional Value-at-Risk，条件风险价值）}，组织成\emph{双层优化}\cite{li2023marc}：上层枚举 ego 策略（沿用 DCP-Tree 思想），下层对每个 ego 策略算 CVaR$_\alpha$ 而非期望——最坏 $(1-\alpha)$ 尾部场景的平均代价。CVaR 用 Rockafellar--Uryasev 形式\cite{rockafellar2000cvar}
\begin{equation}
  \mathrm{CVaR}_\alpha(L)=\min_{t}\Big\{t+\tfrac{1}{1-\alpha}\,\mathbb{E}\big[(L-t)_+\big]\Big\}
  =\min_{t}\Big\{t+\tfrac{1}{1-\alpha}\sum_k P(\omega_k)\,(L_k-t)_+\Big\},
  \label{eq:pr-cvar}
\end{equation}
对离散场景，最优 $t$ 在某个 $L_k$ 处取到，遍历即可。$\alpha$ 是风险旋钮——$\alpha\to 0$ 趋近期望（风险中性）、$\alpha\to 1$ 趋近最坏情况（极度保守）。这样选出的策略对“小概率灾难场景”更稳健（倾向选“即使他车抢行也能安全”的策略，哪怕在“他车让行”的大概率场景下效率略低）。MARC 的运动规划器输出“名义轨迹 + 多条延迟分支”（树形），而非传统的“选中某个场景的单一最优”——保留对所有未来的兼容性直到执行。\pz{控制部深讲：CVaR / 风险敏感 MPC 与分布鲁棒优化（DRO）的完整理论（一致性风险度量公理、DRO ambiguity set）留待控制部，本章只给规划用途的关键式与直觉。}

\begin{insight}
\textbf{为什么偏偏是 CVaR}。其一，\emph{CVaR 是凸的、可优化的}——\cref{eq:pr-cvar} 对决策变量凸（$(L-t)_+$ 凸、取期望与 $\min_t$ 保凸），所以下层是凸优化；对比 VaR（分位数）非凸、不可直接优化。其二，\emph{CVaR 是一致性风险度量（coherent risk measure）}——满足次可加性、正齐次、单调、平移不变四公理，对“分散风险”友好，不会出现 VaR 那种“分散反而显得更危险”的悖论。CVaR 同时满足“对尾部敏感 + 凸可优化 + 一致性”，正好卡在风险敏感规划需要的位置——这也是它几乎成为默认选择的原因。
\end{insight}

\begin{table}[H]\centering\small
\caption{从 EPSILON 到 MARC 的三处关键差异}
\label{tab:pr-marc}
\begin{tabular}{lll}
\toprule
维度 & EPSILON & MARC \\
\midrule
分支点 & 固定（DCP-Tree 每 $2\,$s 节点） & 动态（场景级 divergence，偏差函数 $\mathcal{M}$ 超阈值） \\
代价度量 & 场景概率加权\textbf{期望} & \textbf{CVaR$_\alpha$} 尾部风险 \\
运动规划输出 & 选定策略下的单一最优轨迹 & 名义轨迹 + 多条延迟分支（树形） \\
\bottomrule
\end{tabular}
\end{table}

\begin{insight}
\textbf{MARC = EPSILON 的双层结构 + “期望”换“CVaR” + “固定分支”换“动态分支”}。三处改动都不动 EPSILON 的“行为枚举 + 运动优化”骨架——保骨架、换组件。这也是工程演进的常见形态。两步进化方向不同：EUDM 解决“组合爆炸”（让分支规划可实时，求\emph{效率}），MARC 解决“分支点该放哪 + 怎么对罕见危险稳健”（让分支规划更聪明，求\emph{质量}）。研究有两个永恒维度——更快，或更好。
\end{insight}

\begin{pitfall}[陷阱：把 CVaR 当“最坏情况” / 无脑上动态分支]
其一：以为 CVaR 就是只看最坏那个场景（$\max_k$）——过度保守、退回 frozen 边缘。CVaR$_\alpha$ 是最坏 $(1-\alpha)$ \emph{尾部的平均}，不是单个最坏点；$\alpha=0.9$ 只看最坏 $10\%$ 场景的平均，仍比纯最坏温和。用 \cref{eq:pr-cvar} 实现，验证 $\alpha\to0$ 退化为期望、$\alpha\to1$ 才趋近最坏。其二：以为动态分支点一定比固定的好——在意图揭示时刻本就稳定的场景（如固定路口）上，逐时刻算场景级偏差的开销得不偿失；揭示时刻多变（密集交互）用 MARC 动态分支，揭示时刻稳定用 EUDM 固定分支，这本身就是一次方法选型。
\end{pitfall}

\begin{practice}
\begin{enumerate}
  \item \textbf{[CVaR 选策略]} 在两场景分支里把“期望代价选策略”换成“CVaR$_\alpha$ 选策略”（\cref{eq:pr-cvar}）。构造一个“大概率省时但小概率碰撞”的策略和一个“稍慢但始终安全”的策略，扫 $\alpha$ 观察 MARC 何时从前者切到后者。
  \item \textbf{[动态分支点]} 用 \cref{eq:pr-marc-tbranch} 复现上面的数值示意：给两条会分叉的轨迹，扫阈值 $\theta$，画出动态分支点随 $\theta$ 的变化，对比固定 $2\,$s 分支点。
  \item \textbf{[思考题]} MARC 用场景级偏差函数 $\mathcal{M}$ 判断“场景何时分叉”。若 $\mathcal{M}$ 只看轨迹\emph{均值}之差，而两个场景轨迹均值相同但方差差异大，它能否捕捉这种“分叉”？要捕捉它，应换用何种\emph{分布层面}的度量（如 Wasserstein 距离）？这对分支点选择意味着什么？
\end{enumerate}
\end{practice}

% ════════════════════════════════════════════════════════════════
\section{Multi-stage NMPC 同构与 vs MCTS}
\label{sec:pr-msnmpc}

\subsection{Multi-stage NMPC：化工社区的同一棵树}
\label{sec:pr-msnmpc-iso}

分支规划起源于自动驾驶（Hardy 2013 → MPDM → EUDM）。但在\emph{化工过程控制}社区，Lucia 与 Engell 等独立发展出 \textbf{multi-stage NMPC（多阶段非线性 MPC）}\cite{lucia2013multistage}，用 \emph{scenario tree（场景树）} 处理参数不确定性（如反应速率常数、活化能未知，落在某区间）。把不确定性离散化成几个代表值（最大 / 最小 / 名义），在预测时域上展开成场景树。惊人的是：\textbf{它的数学结构与 contingency planning 严格同构}——

\begin{theorem}[contingency tree 与 scenario tree 的同构]\label{thm:pr-isomorphism}
分支 / contingency 规划（\cref{def:pr-contingency}）与 multi-stage NMPC 是同一个优化骨架的两个实例，对应关系为：场景 $\omega_k$ $\leftrightarrow$ 参数不确定性的离散实现；共享干 + 非预期约束 $\leftrightarrow$ multi-stage 的 here-and-now + 非预期约束；场景加权代价 $\leftrightarrow$ 场景概率加权代价；分支点 $t_{\mathrm{branch}}$ $\leftrightarrow$ 鲁棒时域 $N_r$（robust horizon）。
\end{theorem}

\cref{thm:pr-isomorphism} 不是巧合而是同一数学结构。其中 \textbf{鲁棒时域 $N_r$}：场景树只在前 $N_r$ 步分叉，之后每枝参数固定、不再分叉——否则树随时域指数膨胀 $d^{N_r}$（$d$ = 每步分支因子）。这对应 contingency 里“分支后各枝独立到底”的截断，是控制复杂度的关键旋钮（多数情况 $N_r=1\sim2$ 已足够，与 EUDM“一周期一次切换 + 高频重规划”同理）。multi-stage NMPC 之所以“非保守”，关键在 \emph{recourse（追索 / 反馈在预测中）}：场景树的分支显式编码“未来控制可随不确定性揭示而调整”，于是不必用一条序列同时对所有实现安全——只需共享干段（信息揭示前）对所有实现安全，之后各分支各自应对。这正是分支规划“分支”那一半的作用。

\begin{insight}
\textbf{两个社区爬到了同一个数学山顶却没看见对方}——不是思想不通，而是“语言 + 指标 + 场所”的壁垒：自驾说“他车意图、语义策略”，化工说“参数不确定性、recourse”（关键词不重叠 → 检索不到）；自驾重实时性 / 可解释，化工重最优性 / 递归可行性证明（在意的不同）；自驾发 ICRA / RA-L / T-RO，化工发 J.\ Process Control / Automatica（读的刊不同）。\textbf{启发}：遇到一个结构，先用\emph{数学语言}而非\emph{领域术语}去检索（搜 “scenario tree optimization” “non-anticipativity” “multi-stage stochastic programming”），常能省下重新发明的几年；并能借用化工社区几十年的递归可行性 / 鲁棒性理论来做自驾 contingency（如用 \texttt{do-mpc} / acados 多阶段 OCP 实现 contingency）。这是科研里反复上演的故事（卡尔曼滤波在控制与信号处理、EM 算法在统计与机器学习……）。
\end{insight}

\subsection{分支规划 vs MCTS：与搜索式规划的结构对比}
\label{sec:pr-mcts}

分支规划展开一棵“策略树”，MCTS（蒙特卡洛树搜索，AlphaGo / MuZero 的核心）也展开一棵树。看似都是“树搜索”，但结构、机制、保证截然不同。

\begin{table}[H]\centering\small
\caption{分支规划与 MCTS 的结构对比}
\label{tab:pr-mcts}
\begin{tabular}{p{0.17\textwidth} p{0.37\textwidth} p{0.37\textwidth}}
\toprule
维度 & 分支规划（MPDM / EUDM） & MCTS（AlphaGo / MuZero） \\
\midrule
树的展开 & 领域先验\textbf{固定宽度}策略树（宽度 = 策略数 / 关键交互，深度浅） & \textbf{UCB 自适应}探索-利用（按 $Q+c\sqrt{\ln N/n}$ 选择性深入） \\
Rollout & 闭环仿真（IDM / MOBIL 等\textbf{确定性}反应模型，一次即可） & 随机 rollout 或价值网络估值（需多次平均） \\
反向传播 & \textbf{无}——每周期从零重构树 & \textbf{有}——累积统计量 $N(s,a),Q(s,a)$ 反传回根 \\
可解释性 & 高——每个节点是语义明确的策略 & 低——隐式策略 / 值网络 \\
连续动作 & 天然支持（MPC 风格连续控制） & 需离散化或 Progressive Widening \\
收敛保证 & 无（启发式 + 领域先验） & 有（UCB → 渐近最优 / regret $O(\sqrt{N\ln N})$） \\
实时性 & 强（固定小树，$50$--$100\,$ms） & 取决于模拟次数预算，可重 \\
\bottomrule
\end{tabular}
\end{table}

\begin{insight}
\textbf{分支规划用“领域先验”换“搜索”，MCTS 用“搜索”换“通用性”}。分支规划把“该考虑哪些未来”的知识硬编码进固定宽度策略树，不需大量搜索就能实时决策——代价是依赖人工设计的策略集、对先验外的情况无能为力（漏了“急加速抢过”这个动作就永远做不出）。MCTS 不依赖领域先验，靠 UCB 在巨大搜索空间自适应探索-利用，理论上渐近最优、通用——代价是要大量模拟、黑箱、连续动作难处理。“无反传”不是分支规划偷懒，而是\emph{确定性 rollout 的自然结果}（一次就够，无需累积）；“有反传”也不是 MCTS \rare{啰}嗦，而是\emph{随机 rollout 的必然要求}（单次有方差，须多次累积才收敛）。选哪个取决于你的领域知识有多强、实时预算有多紧、动作空间是离散还是连续。两条路线正在收敛——神经化分支（学习预测分支价值 / 分支点，让 contingency 获得 MCTS 式自适应深度）与 belief-space MCTS（把 AlphaZero 搬进信念空间处理部分可观测）是当前最活跃的中间地带。
\end{insight}

\begin{pitfall}[概念误区：以为分支规划是“简化版 MCTS”]
新手想法：“分支规划不就是没做反向传播、宽度固定的 MCTS 吗？”后果：试图给分支规划加 UCB / 反向传播把它“升级”成 MCTS，结果丢掉实时性和可解释性、又没获得 MCTS 的理论保证。\textbf{根本原因}：两者是\emph{不同的设计取舍}，不是简繁关系。分支规划的固定宽度 + 无反传是\emph{特意为之}（换实时性与可解释），不是 MCTS 的退化。按问题选——强领域先验 + 紧实时预算 + 需可解释 → 分支规划；弱先验 + 大搜索空间 + 容忍黑箱 → MCTS；想要两者之长 → 神经化分支 / belief MCTS。
\end{pitfall}

\begin{practice}
\begin{enumerate}
  \item \textbf{[统一形式化]} 写出一个同时涵盖 MPDM 与 multi-stage NMPC 的统一优化形式（共享干 / 非预期约束 + 分支 + 场景加权代价），分别说明两者是它的哪种特例（场景含义、评估方式、输出类型的取值）。
  \item \textbf{[对比实现]} 在一个小的离散决策场景（如简化路口）里，分别用“固定宽度策略枚举”和“带 UCB 的 MCTS”求解，对比决策质量、计算时间、可解释性。
  \item \textbf{[思考题]} 为什么连续动作对 MCTS 是难点、对分支规划（MPC 风格）却天然？这和两者“如何表示动作空间”有关。
\end{enumerate}
\end{practice}

% ════════════════════════════════════════════════════════════════
\section{Tube MPC 的数学本质}
\label{sec:pr-tube}

至此第一条主线（离散意图的分支规划）告一段落。从本节起转入第二条主线——\emph{连续有界扰动}下的鲁棒规划。Tube MPC 是其中工程生态最成熟、复用价值最高的工具。

\paragraph{动机：名义 MPC 在扰动下为什么不安全。}
名义 MPC（\cref{eq:pr-nominal-mpc}）假设系统严格按 $\mathbf{x}_{k+1}=A\mathbf{x}_k+B\mathbf{u}_k$ 演化。但真实系统有扰动 $\mathbf{w}_k$：每步真实状态被推离预测一点，偏离随时间\emph{累积放大}（若不主动镇定，偏差可能发散）。结果是：名义 MPC 算出的轨迹哪怕“贴着约束边界飞”在理想世界里最优，真实世界里一点扰动就越界。鲁棒规划的提问是：\emph{给定扰动落在已知有界集 $\mathbf{w}_k\in\mathcal{W}$ 内，如何让真实轨迹在所有可能的扰动序列下都满足约束？} Tube MPC 给出的答案优雅得出人意料——它不在线枚举所有扰动（那会爆炸，正是 min-max 方法的困境），而是把不确定闭环系统\emph{分解}成一个确定的“名义系统”加一个被反馈镇定住的“偏差系统”\cite{mayne2005tube}。

\paragraph{两个集合运算（贯穿本节）。}
\begin{definition}[Minkowski 和与 Pontryagin 差]\label{def:pr-minkowski}
\textbf{Minkowski 和} $\mathcal{A}\oplus\mathcal{B}=\{a+b:a\in\mathcal{A},\,b\in\mathcal{B}\}$（把 $\mathcal{B}$ 沿 $\mathcal{A}$ 的每个点扫一遍）；例 $[-1,1]\oplus[-2,2]=[-3,3]$。\textbf{Pontryagin（Minkowski）差} $\mathcal{A}\ominus\mathcal{B}=\{x:x+b\in\mathcal{A}\ \forall b\in\mathcal{B}\}$（把 $\mathcal{A}$ 向内腐蚀 $\mathcal{B}$ 那么多）。它\emph{不是} $\oplus$ 的逆运算：一般 $(\mathcal{A}\ominus\mathcal{B})\oplus\mathcal{B}\subseteq\mathcal{A}$，不取等。
\end{definition}

\subsection{核心推导：分解、偏差动力学、RPI 集、约束收紧}
\label{sec:pr-tube-derivation}

\textbf{第 1 步——问题设定}。带有界加性扰动的线性系统
\begin{equation}
  \mathbf{x}_{k+1}=A\mathbf{x}_k+B\mathbf{u}_k+\mathbf{w}_k,\qquad \mathbf{w}_k\in\mathcal{W},
  \label{eq:pr-tube-sys}
\end{equation}
状态约束 $\mathbf{x}_k\in\mathcal{X}$、控制约束 $\mathbf{u}_k\in\mathcal{U}$，$\mathcal{W}$ 是已知有界凸集（如多面体或盒集），$\mathcal{X},\mathcal{U}$ 凸。目标：设计控制律使真实轨迹对\emph{任意}扰动序列都满足约束。

\textbf{第 2 步——核心分解}。引入不受扰动的\emph{名义系统} $\mathbf{z}_{k+1}=A\mathbf{z}_k+B\mathbf{v}_k$（$\mathbf{z}_k$ 名义状态、$\mathbf{v}_k$ 名义控制，待优化），定义\emph{偏差} $\mathbf{e}_k\triangleq\mathbf{x}_k-\mathbf{z}_k$，采用控制律
\begin{equation}
  \mathbf{u}_k=\mathbf{v}_k+K(\mathbf{x}_k-\mathbf{z}_k)=\mathbf{v}_k+K\mathbf{e}_k,
  \label{eq:pr-tube-law}
\end{equation}
$K$ 是\emph{预先设计好的反馈增益}（离线选定，使 $A+BK$ 稳定，常对名义系统做 LQR 得到，见 \cref{sec:ctrl-lqr}）。名义部分 $\mathbf{v}_k$ 负责“往目标走”，反馈部分 $K\mathbf{e}_k$ 负责“把真实状态拉回名义轨迹”。

\textbf{第 3 步——偏差动力学}。代入三式：
\begin{equation}
\begin{aligned}
\mathbf{e}_{k+1}&=\mathbf{x}_{k+1}-\mathbf{z}_{k+1}
=(A\mathbf{x}_k+B\mathbf{u}_k+\mathbf{w}_k)-(A\mathbf{z}_k+B\mathbf{v}_k)\\
&=A\mathbf{e}_k+B(\mathbf{u}_k-\mathbf{v}_k)+\mathbf{w}_k
=A\mathbf{e}_k+BK\mathbf{e}_k+\mathbf{w}_k
=(A+BK)\,\mathbf{e}_k+\mathbf{w}_k.
\end{aligned}
\label{eq:pr-tube-errdyn}
\end{equation}
\textbf{关键观察}：\cref{eq:pr-tube-errdyn} 是一个\emph{自治的、被 $K$ 镇定的、只受扰动驱动的}线性系统——它和名义轨迹 $\mathbf{v}_k$、和目标都无关。“对抗扰动”被完全隔离到偏差系统里，与“往哪走”（名义系统）解耦。这就是 Tube MPC 把难题拆开的那一刀。

\textbf{第 4 步——RPI 集与不变性}。既然 $A+BK$ 稳定，偏差不发散，那它被困在多大范围里？这是 \textbf{RPI（Robustly Positively Invariant，鲁棒正不变）集}的概念。
\begin{definition}[RPI 集]\label{def:pr-rpi}
集合 $\Omega$ 对偏差动力学 \cref{eq:pr-tube-errdyn} 是 RPI 的，若
\begin{equation}
  (A+BK)\,\Omega\ \oplus\ \mathcal{W}\ \subseteq\ \Omega.
  \label{eq:pr-rpi-cond}
\end{equation}
\end{definition}
\textbf{含义与证明}：假设某时刻 $\mathbf{e}_k\in\Omega$，则 $\mathbf{e}_{k+1}=(A+BK)\mathbf{e}_k+\mathbf{w}_k\in(A+BK)\Omega\oplus\mathcal{W}\subseteq\Omega$，仍在 $\Omega$ 里。由归纳法，只要初始 $\mathbf{e}_0\in\Omega$，则 $\mathbf{e}_k\in\Omega$ 对所有 $k$、对任意扰动序列成立。一句话：\emph{偏差一旦进入 $\Omega$ 就永远留在 $\Omega$ 里}。这个 $\Omega$ 就是“管道的横截面”——真实轨迹 $\mathbf{x}_k=\mathbf{z}_k+\mathbf{e}_k$ 始终落在以名义轨迹 $\mathbf{z}_k$ 为中心、形状为 $\Omega$ 的管道内。

\textbf{第 5 步——最小 RPI 集（mRPI）及其计算}。满足 RPI 条件的 $\Omega$ 有很多（越大越易满足），我们想要\emph{最小}的以减少保守度。理论上的最小 RPI 集为
\begin{equation}
  \Omega_\infty=\bigoplus_{i=0}^{\infty}(A+BK)^i\,\mathcal{W},
  \label{eq:pr-mrpi}
\end{equation}
即把当前与所有历史扰动经闭环动力学传播后的影响全加起来。问题是这是\emph{无穷 Minkowski 和}，一般不可有限精确计算。工程上用 Raković 等\cite{rakovic2005invariant} 给出的方法：计算一个 RPI 的\emph{外逼近} $\Omega_\alpha\supseteq\Omega_\infty$，它可被任意精确地逼近 $\Omega_\infty$，且本身是 RPI 集（这套 $\varepsilon$-mRPI 算法是 Tube MPC 离线阶段的标准步骤）。

\textbf{第 6 步——约束收紧}。把“真实轨迹满足约束”翻译成“名义轨迹满足什么约束”。由 $\mathbf{x}_k=\mathbf{z}_k+\mathbf{e}_k$ 且 $\mathbf{e}_k$ 可取遍 $\Omega$：要保证 $\mathbf{x}_k\in\mathcal{X}$ 对所有 $\mathbf{e}_k\in\Omega$ 成立，需 $\mathbf{z}_k+e\in\mathcal{X}\ \forall e\in\Omega$，即
\begin{equation}
  \mathbf{z}_k\in\mathcal{X}\ominus\Omega;
  \qquad\text{同理}\qquad
  \mathbf{v}_k\in\mathcal{U}\ominus K\Omega.
  \label{eq:pr-tube-tighten}
\end{equation}
多面体约束 $C\mathbf{x}\le d$ 收紧为 $C\mathbf{z}\le d-h_\Omega(C^\top)$，其中 $h_\Omega$ 是 $\Omega$ 的\emph{支撑函数}（逐行减去 $\Omega$ 在该方向上的支撑函数值，即 $\Omega$ 在约束法向上的“宽度”，而非 $\Omega$ 的“半径”）。

\textbf{第 7 步——在线只解名义 MPC}。综合起来，Tube MPC 在线求解的是一个\emph{和名义 MPC 同结构、只是约束被收紧}的优化问题：
\begin{equation}
  \min_{\mathbf{v}_{0:N-1}}\ \sum_{k=0}^{N-1}\ell(\mathbf{z}_k,\mathbf{v}_k)+\ell_f(\mathbf{z}_N)\quad\text{s.t.}\quad \mathbf{z}_{k+1}=A\mathbf{z}_k+B\mathbf{v}_k,\ \mathbf{z}_k\in\mathcal{X}\ominus\Omega,\ \mathbf{v}_k\in\mathcal{U}\ominus K\Omega,
  \label{eq:pr-tube-online}
\end{equation}
初始条件通常取 $\mathbf{z}_0$ 使 $\mathbf{x}_0-\mathbf{z}_0\in\Omega$（如 $\mathbf{z}_0=\mathbf{x}_0$，或把 $\mathbf{z}_0$ 设为决策变量约束 $\mathbf{x}_0-\mathbf{z}_0\in\Omega$ 以真正激活反馈 $K$）。解出 $\mathbf{v}_0$ 后，实际施加 $\mathbf{u}_0=\mathbf{v}_0+K(\mathbf{x}_0-\mathbf{z}_0)$。\textbf{离线一次性算 $K$ 和 $\Omega$，在线只解这个收紧的名义 QP}。

\begin{insight}
\textbf{Tube MPC 把“在线对抗扰动”这件难事，拆成了“离线准备”和“在线求解名义问题”两块}。离线：算反馈增益 $K$ 和不变管道 $\Omega$（一次性、可慢慢算）；在线：只解一个和名义 MPC 同样大小的优化（快）。鲁棒性的代价主要由离线买单，在线计算几乎不增加——这正是它能上嵌入式实时系统的根本原因。而约束收紧 $\mathcal{X}\ominus\Omega$ 把“保守度”显式量化成一个集合运算：$\Omega$ 越大（扰动越强或镇定越弱），可行域被腐蚀越多、越保守。你不是凭感觉加裕度，而是让裕度等于 $\Omega$——一个由 $\mathcal{W}$ 和 $K$ 算出来的集合。
\end{insight}

\begin{derivation}[标量例：黄金比的 RPI 集]\label{der:pr-tube-scalar}
取标量系统 $x_{k+1}=x_k+u_k+w_k$（$A=B=1$），$|w|\le0.1$，约束 $|x|\le1,|u|\le0.5$，$Q=R=1$。\textbf{LQR 增益}：离散 Riccati $P=1+P-\frac{P^2}{1+P}$ 化简为 $P^2-P-1=0$，故 $P=\frac{1+\sqrt5}{2}\approx1.618$（黄金比），增益 $K=-\frac{P}{1+P}\approx-0.618$，闭环 $A+BK\approx0.382$。\textbf{RPI 集}：标量下 mRPI 精确求和——$\Omega=\bigoplus_{i\ge0}(0.382)^i[-0.1,0.1]=\big[-\frac{0.1}{1-0.382},\frac{0.1}{1-0.382}\big]\approx[-0.162,0.162]$（等比级数）。\textbf{约束收紧}：状态 $z\in\mathcal{X}\ominus\Omega=[-0.838,0.838]$，控制 $v\in\mathcal{U}\ominus K\Omega=[-0.4,0.4]$（因 $|K|\cdot0.162\approx0.1$）。约 $16\%$ 的状态范围、$20\%$ 的控制权限被划给鲁棒裕度——这就是为应对 $|w|\le0.1$ 所付的保守度代价。扰动减半则 $\Omega$ 减半；用更强的 $K$（更小的 $A+BK$）则级数收敛更快、$\Omega$ 更小。
\end{derivation}

\begin{pitfall}[陷阱：忘了收紧约束 / 在线优化 $K$ / 死磕最小 RPI]
其一（编程）：把名义 MPC 约束直接当真实约束（在 OCP 里写 $\mathbf{z}_k\in\mathcal{X}$、$\mathbf{v}_k\in\mathcal{U}$），忘了收紧——仿真加扰动后机器人偶尔越界但优化器从不报 infeasible（它解的是没收紧的名义问题，自以为可行）。约束收紧是 Tube MPC 鲁棒保证的\emph{唯一}来源，漏掉它就退化成普通名义 MPC。检验：蒙特卡洛跑 1000 条扰动轨迹，统计实际状态是否始终在 $\mathcal{X}$ 内。其二（概念）：以为 $K$ 和 $\mathbf{v}_k$ 一起在线优化——$K$ 必须\emph{离线固定}，正因 $K$ 固定偏差动力学 $(A+BK)$ 才确定、$\Omega$ 才能离线算；在线优化 $K$ 会破坏“离线准备、在线轻量”的整个架构。其三（思维）：一味追求最小 RPI——$\Omega_\infty$ 不可有限表示，逼近到很高精度后边际收益极小；想大幅减小 $\Omega$ 应去改 $K$（更强镇定）而非死磕逼近精度。
\end{pitfall}

\subsection{三代 Tube 的选型}
\label{sec:pr-tube-gen}

\paragraph{动机：Rigid Tube 为什么常过度保守。}
\cref{sec:pr-tube-derivation} 的 Rigid Tube 用一个\emph{固定}的管道横截面 $\Omega$ 走完全程。但真实系统里“需要多大余量”是时变的——离障碍远时不需要那么宽的管道，机动剧烈时才需要。真实偏差从 $\mathbf{e}_0=0$ 被扰动\emph{逐步撑大}（\cref{der:pr-tube-scalar} 里前几步只累积到 $0.1,0.138,0.153,\dots$，渐增到稳态 $0.162$），Rigid 却每步都按稳态收紧——在偏差还没长起来的前几步白白浪费可行域。三代 Tube 的演化史，本质是一部\emph{逐步让管道“该宽则宽、该窄则窄”以减少保守度的历史}——代价是在线计算逐步增加。

\begin{table}[H]\centering\small
\caption{三代 Tube 选型对比}
\label{tab:pr-tube}
\begin{tabular}{p{0.16\textwidth} p{0.24\textwidth} p{0.24\textwidth} p{0.24\textwidth}}
\toprule
维度 & Rigid (2005) & Homothetic (2012) & Elastic (2016) \\
\midrule
管道横截面 & 固定 RPI 集 $\Omega$ & 相似缩放 $\alpha_k\Omega$ & 参数化可变形多面体 \\
在线决策变量 & 仅 $\mathbf{v}_k$ & $\mathbf{v}_k$ + 缩放 $\alpha_k$ & $\mathbf{v}_k$ + 多面体形状参数 \\
保守度 & 最高 & 中 & 最低 \\
在线计算 & 最低（名义 QP） & 中（名义 QP + 标量） & 高（QP 规模增大） \\
工程推荐 & \textbf{嵌入式 / 实时首选} & 控制能力有余量时升级 & 学术最优性研究为主 \\
\bottomrule
\end{tabular}
\end{table}

三代分别为：\textbf{Rigid Tube}\cite{mayne2005tube}（前身为 Langson 等\cite{langson2004robust} 的 Tube 思想）——管道横截面固定全程不变，最快最简但最保守；\textbf{Homothetic Tube}\cite{rakovic2012homothetic}——管道是固定形状 $\Omega$ 的\emph{几何相似缩放} $\alpha_k\Omega$，缩放因子 $\alpha_k\ge0$ 随时间在线优化（早期 $\alpha_k$ 小、收紧更少、可行域更大），在线多了 $\alpha_k$ 这组标量决策变量；\textbf{Elastic Tube}\cite{rakovic2016elastic}——管道是\emph{参数化可变形的多面体}（不仅缩放还能各面方向独立伸缩），保守度最小但 QP 规模最大。

\begin{insight}
\textbf{这是一条“用在线计算换保守度”的连续谱}。三代 Tube 不是“谁取代谁”，而是同一权衡轴上的三个点——你为减少保守度付出的，恰好是在线优化更多的管道参数（从“无”到“一个标量”到“一组形状参数”）。选型的本质就是问自己：平台在线算力够吗？场景对保守度敏感吗？四足、无人机等\textbf{嵌入式实时}场景\emph{首选 Rigid Tube}（计算最快、保证最强、实现最简，绝大多数情况够用），只有当确认 Rigid 的保守度明显损害性能且平台有富余算力时才升级到 Homothetic；Elastic 在线开销大，目前多用于学术研究“保守度能压到多低”。
\end{insight}

\begin{practice}
\begin{enumerate}
  \item \textbf{[推导]} 自己从头推一遍偏差动力学 \cref{eq:pr-tube-errdyn}，说明为什么它与名义控制 $\mathbf{v}_k$ 无关。再解释：若不加反馈（$K=0$），偏差动力学变成什么？在 $A$ 不稳定时会发生什么？
  \item \textbf{[实现]} 对 2D 双积分器系统 \cref{eq:pr-tube-sys}（$\mathcal{W}$ 取小盒集）实现 Rigid Tube MPC：(1) 对名义系统做 LQR 得 $K$；(2) 用 Raković 等\cite{rakovic2005invariant} 方法算 RPI 外逼近 $\Omega$；(3) 收紧约束 \cref{eq:pr-tube-tighten}；(4) 在线解收紧的名义 MPC \cref{eq:pr-tube-online}。蒙特卡洛 1000 次，验证所有实际轨迹落在管道内、不越界。
  \item \textbf{[思考题]} Homothetic Tube 通过在线优化 $\alpha_k$ 减保守度。从“偏差动力学”角度解释：为什么允许管道随时间缩放，能比固定 $\Omega$ 更贴合真实偏差的演化？（提示：真实偏差从 $\mathbf{e}_0$ 出发逐步被 $K$ 镇定，初期可能小于稳态 $\Omega$。）
\end{enumerate}
\end{practice}

% ════════════════════════════════════════════════════════════════
\section{可达安全包络：FaSTrack / RTD / Funnel}
\label{sec:pr-reach}

\paragraph{动机：把“安全证书”从规划器里解耦出来。}
Tube MPC 把鲁棒性\emph{绑进}了 MPC 求解器（约束收紧是优化问题的一部分）。但有时我们想要另一种架构：用任意一个快的规划器（A$^*$、RRT$^*$、甚至简单 MPC）做规划，而\emph{安全保证来自一个独立的、离线算好的“误差包络”}，与具体用哪个规划器无关——换规划器不影响安全证书。可达性方法走这条路线，共同套路是：\emph{离线预计算“规划器 / 控制器跟踪误差的最坏情形包络”，在线用这个包络膨胀障碍（或约束）}。三种代表方法只在“用什么数学工具算这个误差包络”上不同。\pz{控制部深讲：HJ 可达性（Hamilton-Jacobi PDE）、SOS 规划、追逃微分博弈的完整理论留待控制部；本章只给规划用途（解耦式安全包络）的机制与直觉。}

\begin{table}[H]\centering\small
\caption{三种可达性方法对比}
\label{tab:pr-reach}
\begin{tabular}{p{0.18\textwidth} p{0.24\textwidth} p{0.24\textwidth} p{0.24\textwidth}}
\toprule
维度 & FaSTrack & RTD & Funnel Libraries \\
\midrule
误差包络的计算工具 & HJ 可达性（PDE） & 前向可达集 FRS（SOS） & 李雅普诺夫漏斗（SOS） \\
离线产物 & 跟踪误差集 TEB & 参数空间 FRS & 漏斗库 \\
在线做什么 & 膨胀障碍 + 任意规划器 & 在轨迹参数空间选安全参数 & 拼接漏斗序列 \\
维度可扩展性 & 受 HJ 维度诅咒（约 $\le 5$D） & 中（SOS）；zonotope 版 ARMTD 可达 7-DOF & 中（SOS） \\
与规划器的耦合 & 完全解耦 & 与轨迹参数化耦合 & 与机动库耦合 \\
\bottomrule
\end{tabular}
\end{table}

\textbf{FaSTrack}\cite{herbert2017fastrack}：一个低维\emph{规划模型}（规划快）+ 一个高维\emph{跟踪模型}（贴近真实）。离线对二者的\emph{相对动力学}（看成“跟踪器追、规划器 + 扰动逃”的追逃博弈 $\min_a\max_{u_{\mathrm{plan}}}$），用 HJ 可达性求解最大\emph{跟踪误差集 TEB（Tracking Error Bound）}——它回答“无论规划器怎么动、扰动怎么来，跟踪器与规划器的偏差最坏能有多大”，是该相对动力学 HJ 值函数的一个水平集。在线把所有障碍按 TEB \emph{膨胀} → 任意规划器在膨胀空间用简单规划模型规划 → 跟踪器保证真实轨迹落在规划轨迹的 TEB 邻域内 → 不碰真实障碍。\textbf{局限}：grid-based HJ 计算受\emph{维度诅咒}，直接应用一般限于约 5 维以内的相对动力学（系统分解可缓解）。

\textbf{RTD（Reachability-based Trajectory Design）}\cite{kousik2020rtd}：用\emph{前向可达集 FRS} 而非 HJ 网格。离线对一个\emph{轨迹产生模型}（参数化轨迹族），用 SOS 优化计算 FRS（同时界定跟踪误差）；在线把障碍映射到“会导致碰撞的轨迹参数”集合，于是在线问题变成在\emph{轨迹参数空间}里挑一个安全参数（有限非线性约束，可实时）。它强调 \emph{persistent feasibility（持续可行性）}——规定最短规划时长 + 最短传感器视界，保证“新轨迹总能在旧轨迹执行完前算出来”。\textbf{澄清一处常见误述}：RTD 核心用的是 FRS（SOS 计算）；用 zonotope 表示可达集、扩到 7-DOF 机械臂的是其后续工作 \emph{ARMTD}——把 RTD 笼统说成“用 zonotope”是把两者混淆了。

\textbf{Funnel Libraries}\cite{majumdar2017funnel}：离线沿不同机动预计算一个\emph{漏斗（funnel）库}——每个漏斗是随时间演化的不变集，执行对应机动的反馈控制器时状态保证留在漏斗内（即使有界扰动）；漏斗与一个严格稳定的李雅普诺夫函数关联，用 SOS programming 计算（漏斗 $\approx$ 李雅普诺夫水平集的时变包络）。在线把漏斗按顺序\emph{拼接（sequential composition）}成完整运动计划（前一漏斗出口落在后一入口）。

\begin{insight}
\textbf{TEB / FRS / 漏斗 与 Tube MPC 的 RPI 集是同一类对象}——本质上都是“真实轨迹相对名义 / 规划轨迹的偏差被框住的集合”，都是某种\emph{鲁棒不变（或可达）集}。区别只在用法：Tube MPC 把它绑进 MPC 在线约束收紧（\cref{eq:pr-tube-tighten}），可达性方法把它解耦出来离线算、在线膨胀障碍。两者统一在“鲁棒不变 / 可达集”框架下，差异是\emph{工程架构}而非数学本质。TEB 的大小由“跟踪器能力 vs 规划器激进度”的博弈决定，与 Tube 的 $\Omega$ 由“$K$ vs $\mathcal{W}$”决定完全对应。
\end{insight}

\begin{pitfall}[陷阱：以为 TEB 能算任意高维 / 解耦了规划器就能用任意烂规划器]
其一：以为 HJ 很通用、FaSTrack 能处理任意维度——grid-based HJ 计算量随状态维数\emph{指数增长}（维度诅咒），直接应用一般限于约 5 维以内。高维系统改用 RTD（参数空间 + SOS）、系统分解、或学习化 HJ 近似（如 DeepReach 类）。其二：以为“解耦了规划器”就等于“可以用任意烂规划器”——可达性方法解耦的是\emph{安全证书}，不是\emph{规划质量}。TEB 只保证“不碰膨胀后的障碍”，不保证轨迹优或可行；膨胀后障碍变大，差规划器可能在拥挤场景找不到路。在拥挤场景测规划成功率，而不只测碰撞率。
\end{pitfall}

\begin{practice}
\begin{enumerate}
  \item \textbf{[可达性工具链]} 用 \texttt{hj\_reachability}（JAX）对一个简单系统（如 Dubins 车跟踪移动点）计算跟踪误差集 TEB，可视化误差包络形状，观察 TEB 如何随扰动大小变化。
  \item \textbf{[思考题]} FaSTrack 的 TEB 和 Tube MPC 的 RPI 集 $\Omega$（\cref{def:pr-rpi}）在概念上是什么关系？两者能否统一到同一数学框架下？（提示：都是“偏差被框住的鲁棒不变 / 可达集”，差异在于绑进求解器还是解耦。）
\end{enumerate}
\end{practice}

% ════════════════════════════════════════════════════════════════
\section{GP-MPC：数据驱动的残差与机会约束}
\label{sec:pr-gpmpc}

\paragraph{动机：与其假设最坏扰动，不如从数据里学残差。}
Tube、可达性都假设你\emph{知道}扰动集 $\mathcal{W}$ 或模型误差的界，然后对界内最坏情况买单。但很多模型误差是\emph{系统性的、可学习的}（\cref{sec:pr-unc-types} 的模型不确定性）——四旋翼的气动残差、赛车的轮胎-路面交互、机械臂的摩擦。对这类误差，与其用保守的最坏情形界，不如\emph{从数据里学出这个残差}并量化“学得有多准”，据此\emph{自适应}地收紧约束。这就是 GP-MPC：用高斯过程（GP）学残差，把残差的\emph{均值}喂给动力学预测、\emph{方差}喂给机会约束\cite{hewing2020gpmpc}。\emph{为什么是 GP 而非神经网络}：GP 天生给出预测不确定性（方差）——数据多的区域方差小（信任模型），数据少 / 外推的区域方差大（自动保守）。这个“何时信任、何时保守”的信号是 GP-MPC 的灵魂。

\subsection{标准架构}
\label{sec:pr-gpmpc-arch}

GP-MPC 的标准架构\cite{hewing2020gpmpc}：
\begin{equation}
\begin{aligned}
\text{名义模型（物理先验）:}\quad & \mathbf{x}_{k+1}=f_{\mathrm{nom}}(\mathbf{x}_k,\mathbf{u}_k),\\
\text{GP 残差（数据学习）:}\quad & \mathbf{d}_k\sim\mathcal{GP}\big(\mu_d(\mathbf{x}_k,\mathbf{u}_k),\ \Sigma_d(\mathbf{x}_k,\mathbf{u}_k)\big),\\
\text{均值预测:}\quad & \hat{\mathbf{x}}_{k+1}=f_{\mathrm{nom}}(\mathbf{x}_k,\mathbf{u}_k)+\mu_d(\mathbf{x}_k,\mathbf{u}_k),\\
\text{协方差前向传播:}\quad & \Sigma_{k+1}=g(\Sigma_k,\ \Sigma_d),\\
\text{机会约束:}\quad & \Pr\!\big(\mathbf{x}_k\in\mathcal{X}_{\mathrm{safe}}\big)\ge 1-\delta.
\end{aligned}
\label{eq:pr-gpmpc}
\end{equation}
名义模型给物理先验，GP 学它与真实系统的差（残差 $\mathbf{d}_k$）；组合后用 GP 均值 $\mu_d$ 修正预测、用 GP 方差 $\Sigma_d$ 量化“这块区域信不信得过”。状态分布的协方差 $\Sigma_k$ 沿预测时域往前传播（融合 GP 方差与上一步协方差），再用机会约束保证“真实状态落在安全集内的概率 $\ge 1-\delta$”。

\textbf{机会约束如何变成可解的收紧约束}（高斯情形）：对线性约束 $\Pr(\mathbf{c}^\top\mathbf{x}_k\le b)\ge1-\delta$，因 $\mathbf{c}^\top\mathbf{x}_k\sim\mathcal{N}(\mathbf{c}^\top\mu,\mathbf{c}^\top\Sigma\mathbf{c})$，标准化后取单调的 $\Phi^{-1}$，得确定性收紧约束
\begin{equation}
  \mathbf{c}^\top\hat{\mathbf{x}}_k\ \le\ b-\Phi^{-1}(1-\delta)\,\sqrt{\mathbf{c}^\top\Sigma_k\,\mathbf{c}},
  \label{eq:pr-chance-tighten}
\end{equation}
$\Phi^{-1}$ 是标准正态的分位数函数。收紧量 $\Phi^{-1}(1-\delta)\sqrt{\mathbf{c}^\top\Sigma_k\mathbf{c}}$ 由协方差 $\Sigma_k$ 决定——协方差大（不确定性高）则收紧多、保守；协方差小则进取。常用值：$\delta=0.05$ 时 $\Phi^{-1}(0.95)\approx1.645$，$\delta=0.01$ 时 $\approx2.326$。\pz{控制部 / 机会约束规划深讲：\cref{eq:pr-chance-tighten} 的非高斯推广（Chebyshev / Cantelli 集中不等式、场景法）留待相应章节；此处只给高斯情形的关键式。}

\begin{insight}
\textbf{GP-MPC 的精髓不是“用 GP 替换 $f$”，而是“用 GP 的均值修正 $f$、用 GP 的方差驱动保守度”}。学到的不确定性（方差）直接、自动地转化为约束收紧的力度：方差大 → 收紧多 → 保守，方差小 → 进取。相比 Tube MPC 用一个\emph{固定}的扰动集 $\mathcal{W}$，GP-MPC 用的是\emph{状态相关、数据驱动}的不确定性——这是“让保守度由数据决定”的最优雅范例。它站在三条线的交汇处：鲁棒视角（自适应不确定性替代固定扰动集）、学习视角（model-based RL 的特例，学残差 + 在模型上规划）、机会约束视角（残差驱动的 \cref{eq:pr-chance-tighten} 收紧）。
\end{insight}

\subsection{工程要点与实验证据}
\label{sec:pr-gpmpc-eng}

关键工程选择：\textbf{诱导点数量}——完整 GP 推断复杂度 $O(N^3)$（$N$ 为数据点数），在线跑不动，必须用\emph{稀疏 GP}（诱导点 / FITC 近似），工程上常取约 20 个诱导点以保实时；\textbf{协方差传播近似}——UT（unscented transform）/ 线性化 / moment matching，精度与计算各有取舍；\textbf{残差学什么}——通常学“名义模型解释不了的那部分动力学”（气动、轮胎力），而非整个动力学（保留物理先验大幅减少数据需求）。

\textbf{实验证据}：赛车上\cite{kabzan2019gp}用在线更新的 GP 学轮胎-路面残差，在 AMZ Driverless 全尺寸无人方程式赛车（gotthard）上 GP 字典收敛后单圈时间缩短约 $10\%$；四旋翼上\cite{torrente2021gp}用 GP 学气动残差嵌入 MPC，在速度达 $14\,$m/s、加速度超 $4g$ 的高速飞行中，相较 state-of-the-art 线性阻力模型轨迹跟踪误差最多降低 $70\%$。

\begin{pitfall}[陷阱：在线用完整 GP / 以为方差大就是“数据不够” / 只用均值丢方差]
其一（编程）：把所有历史数据点放进 GP 做完整推断，$O(N^3)$ 随时间无界增长拖垮实时——用稀疏 GP（约 20 诱导点）或在线字典（按信息量增删）控制规模。其二（概念）：以为 GP 方差大就一定是“数据不够、多采就行”——GP 方差混合了 epistemic（外推 / 数据少，采集可消减）和 aleatoric（系统固有噪声地板，采集消不掉）两部分（\cref{sec:pr-unc-types} 的认知 vs 偶然）；若方差主要来自 aleatoric，多采数据无用，应接受它并用机会约束容忍。其三（思维）：只用 GP 均值、忽略方差，把均值当新动力学做普通 NMPC——这等于丢掉 GP-MPC 的全部鲁棒价值；在 GP 外推区（方差大但均值不可靠）盲目自信反而比纯名义 MPC 更危险。均值进动力学、方差进机会约束，两者都用。
\end{pitfall}

\begin{practice}
\begin{enumerate}
  \item \textbf{[GP-MPC 入门]} 用 GPy 或 GPflow 学一个 1D 残差（如给单摆动力学加未建模非线性项），把 GP 均值嵌入 CasADi NMPC。对比有无 GP 残差的跟踪误差，观察 GP 方差在数据稀疏处如何增大、如何提供“何时信任 / 保守”的信号。
  \item \textbf{[协方差传播]} 在上面的 GP-MPC 里实现两种协方差传播（线性化 vs UT），对比它们在预测时域末端给出的 $\Sigma_k$ 大小与机会约束收紧量（\cref{eq:pr-chance-tighten}）的差异。
  \item \textbf{[思考题]} 若用神经网络替代 GP 学残差，你如何估计 NN 的不确定性以驱动机会约束？（提示：MC Dropout、Deep Ensemble；以及如何分解 epistemic / aleatoric。）
\end{enumerate}
\end{practice}

% ════════════════════════════════════════════════════════════════
\section{CBF-QP 安全滤波器}
\label{sec:pr-cbf}

\paragraph{动机：在控制输出端加一道实时“护栏”。}
前面几节都是“\emph{事前}在规划里预留余量”（规划层）。CBF（Control Barrier Function，控制屏障函数）安全滤波走互补的另一条路——“\emph{事中}在控制输出端实时纠偏”（控制层）：让任意一个名义控制器（PD、MPC、甚至黑箱 RL 策略）先给出控制 $\mathbf{u}_{\mathrm{nom}}$，再用一个轻量 QP 在输出端\emph{以最小修正}把它拉回安全集。这道护栏特别适合给“\emph{不可证明安全}的控制器”兜底——比如仿真表现好但偶发危险动作的 RL 策略。\pz{控制部深讲：CBF / CLF 的完整理论（前向不变性证明、与李雅普诺夫稳定性的对偶、ISSf 鲁棒 CBF、自适应 CBF）将在控制部另设章；本章侧重其\emph{规划 / 决策用途}——作为安全滤波层给任意控制器兜底，给出关键式与直觉。}

\paragraph{反面：没有安全滤波时会怎样。}
一个只追求性能的名义控制器（如直奔目标的 PD）在障碍附近会毫不犹豫地朝目标输出控制，哪怕那会撞上障碍——它的目标里没有“安全”这一项。CBF 滤波器的作用，就是在 $\mathbf{u}_{\mathrm{nom}}$ 即将导致越界时，\emph{在所有满足安全约束的控制里挑一个最接近 $\mathbf{u}_{\mathrm{nom}}$ 的}，既尊重名义控制器的意图、又不越界。

\subsection{从安全集到 CBF-QP}
\label{sec:pr-cbf-qp}

把“安全”编码为集合 $\mathcal{C}=\{\mathbf{x}:h(\mathbf{x})\ge 0\}$，$h$ 连续可微（如“离障碍至少 $r$”写成 $h(\mathbf{x})=\|\mathbf{x}-\mathbf{x}_{\mathrm{obs}}\|^2-r^2$）。我们要保证状态\emph{永远留在} $\mathcal{C}$ 内（$\mathcal{C}$ 前向不变）。

\begin{definition}[控制屏障函数条件]\label{def:pr-cbf}
$h$ 是系统的 CBF，若存在控制 $\mathbf{u}$ 使
\begin{equation}
  \dot{h}(\mathbf{x},\mathbf{u})\ \ge\ -\alpha\big(h(\mathbf{x})\big),
  \label{eq:pr-cbf-cond}
\end{equation}
其中 $\alpha(\cdot)$ 是扩展类 $\mathcal{K}_\infty$ 函数（常取线性 $\alpha(h)=\gamma h$，$\gamma>0$；须定义在整个 $\mathbb{R}$ 上以覆盖 $h<0$ 的恢复行为）。
\end{definition}
\textbf{直觉}：当 $h$ 接近 0（贴近边界）时约束趋于 $\dot h\ge0$，强制 $h$ 不再减小、把状态挡在边界内；$h$ 很大（远离边界）时约束很松、几乎不干预\cite{ames2017cbf,ames2019cbf}。

\textbf{对控制仿射系统的可解形式}。对 $\dot{\mathbf{x}}=f(\mathbf{x})+g(\mathbf{x})\mathbf{u}$，由链式法则 $\dot h=\nabla h^\top(f+g\mathbf{u})=L_f h(\mathbf{x})+L_g h(\mathbf{x})\,\mathbf{u}$（$L_f h,L_g h$ 为李导数）。CBF 条件 \cref{eq:pr-cbf-cond} 变成对 $\mathbf{u}$ 的\emph{线性}不等式，于是 \textbf{CBF-QP 安全滤波器}是一个凸二次规划：
\begin{equation}
  \mathbf{u}^*=\arg\min_{\mathbf{u}}\ \|\mathbf{u}-\mathbf{u}_{\mathrm{nom}}\|^2\quad\text{s.t.}\quad L_f h(\mathbf{x})+L_g h(\mathbf{x})\,\mathbf{u}\ge-\alpha(h(\mathbf{x})),\ \ \mathbf{u}\in\mathcal{U}.
  \label{eq:pr-cbf-qp}
\end{equation}
它在“最接近名义控制”的前提下强制满足安全约束，\emph{微秒级可解}，适合 $1\,$kHz+ 的高频安全滤波。下面是单积分器（$\dot{\mathbf{x}}=\mathbf{u}$，即 $f=0,g=I$）+ 圆形障碍的最简实现：

\begin{codebox}[Python]{cbf\_qp\_filter（单积分器 + 圆障碍）}
import numpy as np
import cvxpy as cp

def cbf_qp_filter(x, u_nom, x_obs, r, gamma, u_min, u_max):
    # safe set C = {x : h(x) = ||x - x_obs||^2 - r^2 >= 0}
    n = x.shape[0]
    h = float(x @ x - 2 * x @ x_obs + x_obs @ x_obs - r**2)  # ||x-x_obs||^2 - r^2
    grad_h = 2.0 * (x - x_obs)         # grad h
    Lf_h = 0.0                         # f = 0  -> L_f h = 0
    Lg_h = grad_h                      # g = I  -> L_g h = grad h
    u = cp.Variable(n)
    cons = [Lf_h + Lg_h @ u >= -gamma * h,   # CBF: L_f h + L_g h u >= -gamma h
            u >= u_min, u <= u_max]
    cp.Problem(cp.Minimize(cp.sum_squares(u - u_nom)), cons).solve()
    return u.value
\end{codebox}

\begin{derivation}[标量手算：CBF 把撞墙速度平滑刹到零]\label{der:pr-cbf-scalar}
系统 $\dot x=u$，安全约束“别越过墙 $x\le x_{\max}$”编码为 $h(x)=x_{\max}-x\ge0$。因 $\dot h=-\dot x=-u$，CBF 条件 \cref{eq:pr-cbf-cond}（取 $\alpha(h)=\gamma h$）给 $-u\ge-\gamma(x_{\max}-x)$，即 $u\le\gamma(x_{\max}-x)$。取 $x_{\max}=1,x=0.9,\gamma=2,u_{\mathrm{nom}}=1$（名义想以速度 1 冲墙）：约束 $u\le2(1-0.9)=0.2$，$u_{\mathrm{nom}}=1>0.2$ 违反，QP 取最近可行解 $u^*=0.2$（减速）。当 $x\to1$（贴墙），约束 $u\le2(1-x)\to0$，强制停下。$\gamma$ 大 → 允许更晚刹车（约束在远处更松）；$\gamma$ 小 → 更早减速、更保守。
\end{derivation}

\subsection{高阶 CBF：相对度 > 1}
\label{sec:pr-cbf-hocbf}

\cref{eq:pr-cbf-qp} 假设控制 $\mathbf{u}$ 直接出现在 $\dot h$ 里（$L_g h\not\equiv0$）。但若 $h$ 只与位置有关、而控制是加速度（$L_g h\equiv0$，控制要到 $\ddot h$ 才出现），就需要\textbf{高阶 CBF（HOCBF）}——对 $h$ 反复求导直到控制出现，构造一串嵌套的 CBF 条件\cite{xiao2022hocbf}。

\begin{derivation}[双积分器避障的 HOCBF]\label{der:pr-hocbf}
系统 2D 双积分器 $\dot{\mathbf{p}}=\mathbf{v},\dot{\mathbf{v}}=\mathbf{u}$，安全约束 $h(\mathbf{p})=\|\mathbf{p}-\mathbf{p}_o\|^2-r^2\ge0$（相对度 2，一阶 CBF 失效因 $\dot h$ 里无 $\mathbf{u}$）。取线性 class-$\mathcal{K}$（$\alpha_1,\alpha_2>0$），构造 $\psi_0=h$、$\psi_1=\dot\psi_0+\alpha_1\psi_0=\dot h+\alpha_1 h$，要求 $\dot\psi_1+\alpha_2\psi_1\ge0$。逐阶求导（用 $\dot{\mathbf{p}}=\mathbf{v},\dot{\mathbf{v}}=\mathbf{u}$）：$\dot h=2(\mathbf{p}-\mathbf{p}_o)^\top\mathbf{v}$，$\ddot h=2\mathbf{v}^\top\mathbf{v}+2(\mathbf{p}-\mathbf{p}_o)^\top\mathbf{u}$。代入整理成\emph{对 $\mathbf{u}$ 线性}的约束：
\begin{equation}
  2\mathbf{v}^\top\mathbf{v}+2(\mathbf{p}-\mathbf{p}_o)^\top\mathbf{u}+(\alpha_1+\alpha_2)\,2(\mathbf{p}-\mathbf{p}_o)^\top\mathbf{v}+\alpha_1\alpha_2(\|\mathbf{p}-\mathbf{p}_o\|^2-r^2)\ \ge\ 0.
  \label{eq:pr-hocbf}
\end{equation}
HOCBF-QP 即 $\min_{\mathbf{u}}\|\mathbf{u}-\mathbf{u}_{\mathrm{nom}}\|^2$ s.t.\ \cref{eq:pr-hocbf} $+\ \mathbf{u}\in\mathcal{U}$。\textbf{手算走查}：$\mathbf{p}=(1,0),\mathbf{p}_o=(2,0),r=0.5,\mathbf{v}=(1,0)$（正朝障碍冲），$\alpha_1=\alpha_2=4$。$h=0.75$，$\dot h=2(-1,0)\cdot(1,0)=-2$（负值 → 正在接近）。约束：$\big(2+2(-1,0)\cdot\mathbf{u}\big)+8(-2)+16(0.75)\ge0$ → $2-2u_x-16+12\ge0$ → $u_x\le-1$。即必须在 $x$ 方向以至少 1 的速率减速；若名义 $\mathbf{u}_{\mathrm{nom}}=(2,0)$（朝障碍加速），HOCBF-QP 会否决它、改成最接近的可行解 $u_x=-1$（刹车）。
\end{derivation}

\subsection{CLF-CBF 联立与 CBF/MPC 互补}
\label{sec:pr-cbf-mpc}

\textbf{CLF-CBF 联立}。把控制李雅普诺夫函数（CLF，管稳定 / 到达）与 CBF（管安全）放进同一个 QP，安全硬约束优先、稳定尽力而为（常给 CLF 加松弛变量 $\delta$）：
\begin{equation}
  \min_{\mathbf{u},\delta}\ \|\mathbf{u}-\mathbf{u}_{\mathrm{nom}}\|^2+\rho\,\delta^2\quad\text{s.t.}\quad
  \underbrace{\dot V\le-\lambda V+\delta}_{\text{CLF（软，可松弛）}},\ \
  \underbrace{\dot h\ge-\alpha(h)}_{\text{CBF（硬）}},\ \ \mathbf{u}\in\mathcal{U}.
  \label{eq:pr-clf-cbf}
\end{equation}
这是“安全 + 性能”在控制层的统一——安全是硬约束（绝不违反），稳定 / 到达是软目标（尽力而为，违反时罚松弛）。

\textbf{CBF 与 MPC 的互补：层级协作而非替代}。CBF 不是用来取代 MPC 的，两者是\emph{层级协作}——MPC 处理 $10$--$50$ 步长时域的经济 / 性能优化（给方向），CBF 保证每一步的瞬时安全（给护栏）\cite{grandia2021cbf}。Grandia 等\cite{grandia2021cbf}在 ANYmal 四足上给出代表性范例：把 CBF 地形安全约束\emph{同时}嵌入一个低频 kinodynamic MPC 和一个高频逆动力学（WBC）跟踪控制器，在 stepping-stones（落足点离散稀疏）场景下同时实现安全落足与动态稳定。CBF 还可作为\emph{离散时间 CBF（DT-CBF）}加进 MPC 每一步约束：连续条件 \cref{eq:pr-cbf-cond} 的离散版是 $h(\mathbf{x}_{k+1})\ge(1-\gamma)h(\mathbf{x}_k)$（$0<\gamma\le1$）。\pz{典型腿足控制栈外层 MPC 频率较低、内层跟踪 / 全身控制频率高（约 $1\,$kHz 量级），但具体频率因系统而异，并非固定数值；关键在“长时域优化”与“高频瞬时安全”的分工。}

\begin{insight}
\textbf{CBF 是“安全版的李雅普诺夫”}。李雅普诺夫函数用 $\dot V\le0$ 证稳定（把状态吸引到平衡点，\cref{sec:ctrl-lyap-direct}），CBF 用 $\dot h\ge-\alpha(h)$ 证安全（把状态约束在安全集内）。一个“拉向某点”，一个“挡在某集合内”——换个不等号方向，你已熟悉的稳定性论证就成了安全性论证。而 CBF-QP 是“最小侵入式安全”：平时几乎不动 $\mathbf{u}_{\mathrm{nom}}$，只在接近边界时才出手。正因如此，它能给\emph{任意}名义控制器（含黑箱 RL 策略）套一层形式化安全保证，是连接“经典安全”与“学习控制”的关键接口。
\end{insight}

\begin{pitfall}[陷阱：执行器受限时 infeasible / 以为从任意状态都能恢复 / 以为 $\gamma$ 越大越安全]
其一（编程）：当 $\mathbf{u}\in\mathcal{U}$（执行器能力）与 CBF 约束冲突（如刹不住车），QP 返回 infeasible、控制器卡死或输出 NaN——给 CBF 约束加松弛变量 $\delta\ge0$（$L_f h+L_g h\,\mathbf{u}\ge-\alpha(h)-\delta$，目标里重罚 $\delta$），或在设计阶段保证控制不变集。其二（概念）：以为加了 CBF 就“从任意状态都能恢复安全”——CBF 保证的是\emph{前向不变性}（从安全集内出发\emph{留在}安全集内），不保证从集外（$h<0$）恢复到集内（那是 reach-avoid / 可达性问题，需别的机制，\cref{sec:pr-reach}）。确保初始状态在安全集内。其三（思维）：以为 class-$\mathcal{K}$ 增益 $\gamma$ 越大越安全——$\gamma$ 不是“安全程度”旋钮而是“多晚才开始减速”旋钮：$\gamma$ 大 → 允许 $h$ 下降更快 → 状态可更靠近边界才被拦（更激进）；$\gamma$ 小 → 更早干预（更保守）。两种极端都不好，把 $\gamma$ 当“刹车时机”来调。
\end{pitfall}

\begin{practice}
\begin{enumerate}
  \item \textbf{[实现]} 用 \texttt{cbf\_qp\_filter} 实现单积分器 + 圆形障碍的 CBF 安全滤波。让名义 PD 控制器直冲障碍，验证 CBF 修正后的轨迹不碰撞；扫描 $\gamma$，画出“离障碍最近距离 vs $\gamma$”曲线（对照 \cref{der:pr-cbf-scalar}）。
  \item \textbf{[高阶 CBF]} 对相对度为 2 的系统（位置安全约束、加速度作为控制）实现 HOCBF（\cref{eq:pr-hocbf}），对比它与“直接对位置用一阶 CBF”（会失败，因 $L_g h\equiv0$）的区别。
  \item \textbf{[思考题]} 为什么 CBF-QP 能给 RL 策略兜底？它和“在 RL 训练的 reward 里加安全惩罚”有什么本质区别？（提示：部署时的硬保证 vs 训练时的软引导。）
\end{enumerate}
\end{practice}

% ════════════════════════════════════════════════════════════════
\section{两条主线的衔接：嵌套、递归可行性、从 contingency 到博弈}
\label{sec:pr-branch-extra}

本章两条主线（离散意图分支 / 连续扰动鲁棒）并非孤立，这里补三点把它们焊在一起。

\textbf{分支 $\times$ 机会约束的嵌套}。设 $K$ 个离散意图 $\omega_k$（他车让 / 不让），\emph{每个意图内部}障碍位置还带连续高斯不确定 $\Sigma_o^{(k)}$（感知噪声）。联合形式：外层是离散意图分支（本章第一条主线），内层每条分支的避障约束用 \cref{eq:pr-chance-tighten} 的解析高斯收紧处理该意图下的连续不确定：
\begin{equation}
  \min_{\mathbf{u}_{\mathrm{trunk}},\{\mathbf{u}_k\}}\ \sum_{k=1}^{K}P(\omega_k)\,J(\cdot;\omega_k)\ \ \text{s.t.}\ \ \mathbf{u}(t_0{:}t_{\mathrm{branch}})=\mathbf{u}_{\mathrm{trunk}},\ \
  \mathbf{a}_j^\top\bar{\mathbf{x}}_{k,t}\le b_j^{(k)}-\Phi^{-1}(1-\delta)\sqrt{\mathbf{a}_j^\top(\Sigma_{k,t}+\Sigma_o^{(k)})\mathbf{a}_j}\ \forall k,t,j.
  \label{eq:pr-nested}
\end{equation}
\emph{分支回答“哪个未来”，收紧回答“那个未来里多大扰动”}——两者维度正交（意图离散多模态、扰动连续单峰），各管各的，正是 \cref{sec:pr-unc-faces} “两副面孔可叠加”的精确落地。

\textbf{递归可行性与共享干}。分支规划作为滚动重规划的 MPC，也面临\emph{递归可行性（recursive feasibility）}——这一步可行能否保证下一步仍可行。共享干是双刃：它强制各场景共用一段动作，增强短期一致性（治振荡），但也可能损害递归可行性（若把 ego 带到一个“对某新出现的危险场景无法补救”的状态）。缓解手段沿用 Tube / multi-stage 思想：终端约束 / 终端集（要求每条分支末端进入“安全可维持”的鲁棒控制不变集，保证后续可补救）、短共享干 + 高频重规划、软约束兜底（碰撞约束加松弛 + 重罚，保证 QP 永远有解，避免硬 infeasible 卡死）。\pz{控制部 MPC 章深讲：递归可行性与终端集 / 终端代价的稳定性证明（终端代价 + 终端集 = 闭环李雅普诺夫函数）见 \cref{sec:ctrl-mpc-feasibility,sec:ctrl-mpc-stability}；本章只点明共享干带来的新考量。}

\textbf{从 contingency 到博弈}。本章把他车意图当\emph{外生情景}——给定 $\omega_k$ 与概率 $P(\omega_k)$，他车不随 ego 的策略改变其意图分布。但真实强交互里，他车也在\emph{对 ego 最优反应}（你让它就进、你抢它就退），这就从“分支规划”过渡到“博弈规划”（求 Nash / Stackelberg 均衡）。判断用哪个：他车对 ego 的反应\emph{弱}（如 ego 是小车、他车是大流量主路）→ contingency 够用；反应\emph{强}（势均力敌的博弈、谈判式交互）→ 需要博弈。Contingency Games 把两者结合：在博弈框架里做 contingency。\pz{博弈规划（Nash / Stackelberg 均衡、iterative best response、Contingency Games）若本书另设章，可在此 \texttt{cref} 桥接；本章只点明 contingency 是博弈在“他车反应弱”极限下的特例。}

\begin{insight}
\textbf{选规划方法，本质是在“先验”和“搜索 / 数据”之间分配资源}。这条权衡贯穿全章——鲁棒规划用强模型先验（动力学 + 不确定性界）做解析鲁棒 / 收紧，分支规划用强领域先验（语义策略）做枚举分支，而 RL / MCTS 在先验稀薄处用搜索 / 数据补。你对问题了解越多（强先验），越能用轻量结构（固定树、解析约束）实时求解；了解越少，越要靠搜索或数据。\textbf{在安全关键机器人里，“可验证安全”常压倒“平均性能最优”}——这是分支 / Tube / CBF 等经典方法长盛不衰的根本原因：一个平均更好但偶尔莫名失败、无法验证的黑箱在自驾 / 医疗里可能根本无法部署。学习与经典不是替代关系，而是“学习提供能力、经典提供可验证骨架”的互补。
\end{insight}

% ════════════════════════════════════════════════════════════════
\section*{常见误解汇总}
\addcontentsline{toc}{section}{常见误解汇总}

\begin{enumerate}
  \item \textbf{“分支规划 = 为每个场景各跑一个独立 MPC。”} 漏掉了共享干（非预期约束）——分支规划的灵魂是“分支前必须共用一个当前动作”，正是这条耦合约束让“当前该做什么”有唯一稳健答案（\cref{sec:pr-branch-formal}）。
  \item \textbf{“分支越多（$K$ 越大）越安全。”} $K$ 是计算成本的主因，安全主要由“是否覆盖了真正危险的那几个模态”决定（CFB 思想），不由数量决定（\cref{sec:pr-branch-knobs}）。
  \item \textbf{“MPDM 在预测他车的精确轨迹。”} MPDM 预测的是离散策略的概率，轨迹是用 IDM / MOBIL 在闭环前向仿真里跑出来的（\cref{sec:pr-mpdm-algo}）。
  \item \textbf{“分支规划是简化版 MCTS。”} 两者是不同的设计取舍（先验密集 vs 搜索密集），不是简繁关系；分支规划的固定宽度 + 无反传是特意为之（\cref{sec:pr-mcts}）。
  \item \textbf{“把名义 MPC 约束直接当真实约束就行。”} 忘了约束收紧——它是 Tube MPC 鲁棒保证的唯一来源，漏掉就退化成普通名义 MPC（\cref{sec:pr-tube-derivation}）。
  \item \textbf{“更先进的 Elastic Tube 总是更好的工程选择。”} 保守度低有代价（在线 QP 规模大）；嵌入式平台往往实时性优先，默认 Rigid（\cref{sec:pr-tube-gen}）。
  \item \textbf{“FaSTrack 的 TEB 能算任意高维系统。”} grid-based HJ 受维度诅咒，直接应用一般限于约 5 维以内（\cref{sec:pr-reach}）。
  \item \textbf{“GP 方差大就是数据不够，多采就行。”} 方差混合 epistemic（可消减）与 aleatoric（消不掉的噪声地板）；若 aleatoric 主导，多采无用（\cref{sec:pr-gpmpc-eng}）。
  \item \textbf{“CBF 的 class-$\mathcal{K}$ 增益 $\gamma$ 越大越安全。”} $\gamma$ 是“刹车时机”旋钮不是“安全程度”旋钮——大则更晚介入（更激进），两种极端都不好（\cref{sec:pr-cbf-mpc}）。
  \item \textbf{“CVaR 就是最坏情况。”} CVaR$_\alpha$ 是最坏 $(1-\alpha)$ 尾部的\emph{平均}，不是单个最坏点；$\alpha\to0$ 退化为期望（\cref{sec:pr-marc-cvar}）。
\end{enumerate}

% ════════════════════════════════════════════════════════════════
\section*{本章小结}
\addcontentsline{toc}{section}{本章小结}

\subsection*{符号表}
\begin{table}[H]\centering\small
\begin{tabular}{ll}
\toprule
符号 & 含义 \\
\midrule
$\omega_k,\ P(\omega_k)$ & 离散未来情景 / 意图及其概率（\cref{eq:pr-contingency}） \\
$\mathbf{u}_{\mathrm{trunk}},\ \{\mathbf{u}_k\}$ & 共享干控制 / 各分支控制 \\
$t_{\mathrm{branch}}$ & 分支点（共享干长度，信息揭示时刻） \\
$\Pi_{\mathrm{ego}},\ \Pi_{\mathrm{other}}$ & ego / 他车的语义策略集 \\
$n,\ n_{\mathrm{key}}$ & 相关他车数 / 关键交互车数（CFB） \\
$\mathcal{X},\mathcal{U},\mathcal{W}$ & 状态约束集 / 控制约束集 / 扰动集 \\
$\mathcal{A}\oplus\mathcal{B},\ \mathcal{A}\ominus\mathcal{B}$ & Minkowski 和 / Pontryagin 差（\cref{def:pr-minkowski}） \\
$\mathbf{z}_k,\mathbf{v}_k$ & 名义状态 / 名义控制 \\
$\mathbf{e}_k=\mathbf{x}_k-\mathbf{z}_k$ & 偏差（真实 $-$ 名义） \\
$K$ & 离线镇定反馈增益（使 $A+BK$ 稳定） \\
$\Omega,\ \Omega_\infty$ & RPI 集 / 最小 RPI 集（\cref{def:pr-rpi,eq:pr-mrpi}） \\
$\alpha_k$ & Homothetic Tube 管道缩放标量序列 \\
$\mu_d,\Sigma_d$ & GP 残差的均值 / 方差（\cref{eq:pr-gpmpc}） \\
$\Sigma_k$ & 状态分布协方差（沿时域传播） \\
$\Phi^{-1}(1-\delta)$ & 标准正态分位数（机会约束收紧系数） \\
$h(\mathbf{x}),\ \mathcal{C}$ & CBF 屏障函数 / 安全集 $\{h\ge0\}$ \\
$\alpha(\cdot),\ \gamma$ & CBF 的扩展类 $\mathcal{K}$ 函数 / 线性增益 \\
$L_f h,L_g h$ & $h$ 沿 $f,g$ 的李导数 \\
$\mathrm{CVaR}_\alpha,\ \alpha$ & 条件风险价值 / 风险置信水平（\cref{eq:pr-cvar}） \\
$N_r$ & multi-stage NMPC 的鲁棒时域 \\
\bottomrule
\end{tabular}
\end{table}

\subsection*{定理与关键结论速查}
\begin{table}[H]\centering\small
\begin{tabular}{p{0.30\textwidth}p{0.62\textwidth}}
\toprule
结论 & 速记 \\
\midrule
分支规划形式化（\cref{def:pr-contingency}） & 场景加权代价 + 共享干（非预期约束）+ 各分支独立；只施加共享干首控制后滚动重规划 \\
MPDM 复杂度（\cref{eq:pr-mpdm-complexity}） & $O(|\Pi_{\mathrm{ego}}|\cdot|\Pi_{\mathrm{other}}|^n)$，随他车数指数爆炸 \\
EUDM 复杂度（\cref{eq:pr-eudm-complexity}） & DCP-Tree 砍底数幂、CFB 砍指数 → $O(d|\Pi_{\mathrm{ego}}|\cdot|\Pi_{\mathrm{other}}|^{n_{\mathrm{key}}})$，$50$--$100\,$ms 可解 \\
同构（\cref{thm:pr-isomorphism}） & contingency tree = scenario tree；共享干 = 非预期、分支点 = 鲁棒时域 $N_r$ \\
Tube 偏差动力学（\cref{eq:pr-tube-errdyn}） & $\mathbf{e}_{k+1}=(A+BK)\mathbf{e}_k+\mathbf{w}_k$，与名义控制无关 → 解耦 \\
RPI 不变性（\cref{eq:pr-rpi-cond}） & $(A+BK)\Omega\oplus\mathcal{W}\subseteq\Omega$ → 偏差进 $\Omega$ 永留 $\Omega$ \\
约束收紧（\cref{eq:pr-tube-tighten}） & $\mathbf{z}_k\in\mathcal{X}\ominus\Omega,\ \mathbf{v}_k\in\mathcal{U}\ominus K\Omega$；多面体减支撑函数 \\
机会约束收紧（\cref{eq:pr-chance-tighten}） & $\mathbf{c}^\top\hat{\mathbf{x}}\le b-\Phi^{-1}(1-\delta)\sqrt{\mathbf{c}^\top\Sigma\mathbf{c}}$ \\
CBF 条件（\cref{eq:pr-cbf-cond}） & $\dot h\ge-\alpha(h)$ → 前向不变；安全版李雅普诺夫 \\
CBF-QP（\cref{eq:pr-cbf-qp}） & 最小修正 $\min\|\mathbf{u}-\mathbf{u}_{\mathrm{nom}}\|^2$ s.t.\ 线性 CBF 约束，微秒级 \\
CVaR（\cref{eq:pr-cvar}） & 最坏 $(1-\alpha)$ 尾部平均；R-U 形式凸、一致性风险度量 \\
\bottomrule
\end{tabular}
\end{table}

\subsection*{知识点总表（含难度）}
\begin{table}[H]\centering\small
\begin{tabular}{p{0.32\textwidth}p{0.42\textwidth}c}
\toprule
知识点 & 一句话掌握 & 难度 \\
\midrule
不确定性四类 / 两副面孔 & 离散意图用分支、连续扰动用鲁棒，正交可叠加 & $\bigstar$ \\
分支规划形式化 & 共享干 + 加权分支 + 非预期约束 & $\bigstar\bigstar$ \\
三个设计旋钮 & 分支点最关键（太早振荡、太晚保守） & $\bigstar\bigstar$ \\
MPDM & 策略枚举 + 闭环前向仿真（IDM/MOBIL） & $\bigstar\bigstar\bigstar$ \\
EUDM & DCP-Tree（砍底数）+ CFB（砍指数）压回实时 & $\bigstar\bigstar\bigstar$ \\
EPSILON 分层 & 行为层定通道 + SSC 凸走廊 QP 填轨迹 & $\bigstar\bigstar\bigstar$ \\
MARC & 动态分支点（场景级 divergence）+ CVaR 风险 & $\bigstar\bigstar\bigstar\bigstar$ \\
multi-stage NMPC 同构 & contingency tree = scenario tree & $\bigstar\bigstar\bigstar$ \\
分支 vs MCTS & 先验密集 vs 搜索密集 & $\bigstar\bigstar$ \\
Tube MPC 本质 & 名义系统 + 反馈 + RPI 集 + 约束收紧 & $\bigstar\bigstar\bigstar$ \\
三代 Tube 选型 & 用在线计算换保守度的连续谱 & $\bigstar\bigstar$ \\
可达包络 & TEB/FRS/漏斗 = 解耦式 RPI 集 & $\bigstar\bigstar\bigstar$ \\
GP-MPC & 均值修正动力学 + 方差驱动机会约束 & $\bigstar\bigstar\bigstar$ \\
CBF-QP 安全滤波 & 最小修正强制前向不变；高阶相对度用 HOCBF & $\bigstar\bigstar\bigstar$ \\
CLF-CBF 联立 & 安全硬约束 + 稳定软目标（松弛） & $\bigstar\bigstar\bigstar$ \\
\bottomrule
\end{tabular}
\end{table}

\subsection*{方法选型决策指南}
\begin{table}[H]\centering\small
\begin{tabular}{p{0.40\textwidth}p{0.52\textwidth}}
\toprule
你的问题 & 推荐方法 \\
\midrule
他人意图离散多模态（让 / 不让、左 / 右） & 分支规划：弱交互用 MPDM/EUDM/EPSILON，强交互用博弈 \\
意图揭示时刻多变 + 怕罕见灾难 & MARC（动态分支点 + CVaR） \\
连续有界扰动、要 100\% 约束满足、嵌入式 & Rigid Tube MPC（首选）；算力有余且保守度是瓶颈才升 Homothetic \\
要换规划器不影响安全证书 & 可达包络（低维 FaSTrack、高维 RTD、机动库 Funnel） \\
模型有系统性可学残差（气动 / 轮胎 / 摩擦） & GP-MPC（均值修正 + 方差收紧） \\
给任意 / 黑箱控制器兜瞬时安全底线 & CBF-QP 安全滤波（相对度 > 1 用 HOCBF） \\
真实安全系统（综合） & 分支处理意图 + 每分支机会约束处理扰动 + CBF 兜底 \\
\bottomrule
\end{tabular}
\end{table}

% ════════════════════════════════════════════════════════════════
\section*{延伸阅读}
\addcontentsline{toc}{section}{延伸阅读}

\textbf{分支 / 场景规划}：概念奠基 Hardy--Campbell\cite{hardy2013contingency}（首次在概率障碍预测上同时优化多条 contingency 路径）；离散策略枚举 MPDM\cite{cunningham2015mpdm}；引导分支压复杂度 EUDM\cite{zhang2020eudm} → 完整系统 EPSILON\cite{ding2021epsilon}（含 SSC 运动层\cite{ding2019ssc}）；动态分支 + 风险 MARC\cite{li2023marc}；化工社区同构的 multi-stage NMPC\cite{lucia2013multistage}。这条线的演进可概括为“先把思想做对（Hardy）、再做快做全（EUDM/EPSILON）、最后做聪明（MARC）”。
\textbf{鲁棒 Tube MPC}：Tube 思想前身\cite{langson2004robust}，Rigid Tube\cite{mayne2005tube}，RPI 集计算\cite{rakovic2005invariant}，Homothetic\cite{rakovic2012homothetic} 与 Elastic\cite{rakovic2016elastic}。
\textbf{可达安全包络}：FaSTrack\cite{herbert2017fastrack}、RTD\cite{kousik2020rtd}、Funnel Libraries\cite{majumdar2017funnel}。
\textbf{数据驱动}：GP-MPC 架构\cite{hewing2020gpmpc}、赛车应用\cite{kabzan2019gp}、四旋翼应用\cite{torrente2021gp}。
\textbf{CBF 安全滤波}：CBF 理论\cite{ames2017cbf,ames2019cbf}、高阶 CBF\cite{xiao2022hocbf}、CBF-MPC/WBC 互补\cite{grandia2021cbf}、风险度量 CVaR\cite{rockafellar2000cvar}、微观交通模型 IDM\cite{treiber2000idm} 与 MOBIL\cite{kesting2007mobil}。

% ════════════════════════════════════════════════════════════════
\section*{与后续 / 相关章节的关系}
\addcontentsline{toc}{section}{与后续 / 相关章节的关系}

\begin{itemize}
  \item \textbf{上承} \cref{ch:planning}（规划导论）：本章把规划三范式（搜索 / 采样 / 优化）扩展到“未来不确定”的情形——单轨迹升级为轨迹树（分支）或轨迹 + 管道（鲁棒）。
  \item \textbf{紧靠} \cref{ch:control}（控制导论）：Tube MPC 的反馈增益来自 LQR（\cref{sec:ctrl-lqr}）、在线问题是收紧的 MPC（\cref{sec:ctrl-mpc}）；CBF 是李雅普诺夫论证（\cref{sec:ctrl-lyapunov}）换不等号方向；递归可行性 / 终端集见 \cref{sec:ctrl-mpc-feasibility,sec:ctrl-mpc-stability}。
  \item \textbf{控制部深讲（本章按需简述，留待深推）}：CBF / CLF 的完整前向不变性与 ISSf 鲁棒理论、HJ 可达性 PDE 与追逃博弈、CVaR / 风险敏感 MPC 与分布鲁棒优化（DRO）——这些理论本章只给\emph{规划 / 决策用途}的关键式与直觉，最深的推导留待控制部相应章节，避免跨部重复。
  \item \textbf{可衔接的专题}：POMDP / belief-space 规划（EUDM 的形式化背景、belief MCTS）、博弈规划（contingency 的强交互推广）、机会约束规划（\cref{eq:pr-chance-tighten} 的系统化展开）——若本书后续设这些章，可在本章对应 \texttt{pz} 标注处补 \texttt{cref} 桥接。
\end{itemize}

% ════════════════════════════════════════════════════════════════
\section*{故障排查手册}
\addcontentsline{toc}{section}{故障排查手册}

\begin{table}[H]\centering\small
\begin{tabular}{p{0.30\textwidth}p{0.30\textwidth}p{0.32\textwidth}}
\toprule
症状 & 可能原因 & 排查 / 修复 \\
\midrule
机器人在人群 / 车流里僵住不动 & 单轨迹规划对悲观预测无解（frozen） & 改分支规划，为每个离散未来留一条出路、当前只走共享干（\cref{sec:pr-branch-formal}） \\
决策在加速 / 刹车间反复横跳 & 缺共享干，最优解随预测翻转突变 & 加非预期约束（约束在控制上）；把 $t_{\mathrm{branch}}$ 设大检验是否靠共享干稳住 \\
分支规划求解超时 & $K$ 太大 / 枚举他车全组合 & 用 CFB 只对关键交互车分支（\cref{sec:pr-eudm-cfb}），$K$ 控制在 $3$--$8$ \\
Tube MPC 仿真加扰动后偶尔越界、但从不报 infeasible & 约束没收紧，解的是名义问题 & 收紧为 $\mathcal{X}\ominus\Omega,\mathcal{U}\ominus K\Omega$（\cref{eq:pr-tube-tighten}）；蒙特卡洛 1000 条核验 \\
Tube MPC 过度保守、绕远路 & $\Omega$ 太大（扰动悲观或 $K$ 太弱） & 重新设计更强的 $K$（更小 $A+BK$）；或升级 Homothetic（\cref{sec:pr-tube-gen}） \\
FaSTrack 算 TEB 内存 / 时间爆炸 & grid-based HJ 维度诅咒 & 相对动力学 $>5$D 改用系统分解 / RTD / 学习化 HJ（\cref{sec:pr-reach}） \\
GP-MPC 每步越来越慢 & 在线用完整 GP，$O(N^3)$ & 用稀疏 GP（约 20 诱导点）/ 在线字典 \\
GP-MPC 在外推区翻车 & 只用均值、丢了方差 & 方差进机会约束（\cref{eq:pr-chance-tighten}）；外推区方差应增大、自动保守 \\
CBF-QP 返回 infeasible / NaN & 执行器 $\mathcal{U}$ 与 CBF 约束冲突 & 加松弛 $\delta\ge0$ 并重罚；或设计控制不变集 \\
CBF 加了仍越界 & 从安全集外（$h<0$）启动 & CBF 只保前向不变；确保初始 $h\ge0$，集外恢复用可达性方法 \\
位置安全约束、加速度控制，CBF 不起作用 & 相对度 > 1，$L_g h\equiv0$ & 用 HOCBF 反复求导直到控制出现（\cref{eq:pr-hocbf}） \\
长时域跑久了闭环发散 & 缺终端集 / 递归可行性 & 加终端集 / 终端代价（\cref{sec:ctrl-mpc-stability}）；测长时域闭环而非单步 \\
\bottomrule
\end{tabular}
\end{table}
